% petz_recovery_supplemental.tex
% Supplemental Material for:
% "The Arrow of Time as Petz Recovery Failure"
% Compile with: pdflatex petz_recovery_supplemental.tex

\documentclass[prl,twocolumn,superscriptaddress,longbibliography]{revtex4-2}

\usepackage{amsmath}
\usepackage{amssymb}
\usepackage{amsfonts}
\usepackage{amsthm}
\usepackage{bm}
\usepackage{hyperref}
\usepackage{booktabs}

\newtheorem{theorem}{Theorem}

% Convenience macros (same as main text)
\newcommand{\kB}{k_{\mathrm{B}}}
\newcommand{\Tr}{\mathrm{Tr}}
\newcommand{\id}{\mathrm{id}}
\newcommand{\Petz}{\mathcal{R}}
\newcommand{\chan}{\mathcal{N}}
\newcommand{\hilb}{\mathcal{H}}
\newcommand{\code}{\mathcal{C}}
\newcommand{\KL}{D}
\newcommand{\ket}[1]{|#1\rangle}
\newcommand{\bra}[1]{\langle#1|}
\newcommand{\braket}[2]{\langle#1|#2\rangle}
\newcommand{\op}[2]{|#1\rangle\langle#2|}
\newcommand{\abs}[1]{\left|#1\right|}
\newcommand{\norm}[1]{\left\|#1\right\|}
\newcommand{\tauasym}{\tau}

\begin{document}

\title{Supplemental Material:\\
Universal Recovery: The Petz Map as Retrodiction Functor,\\
Error Corrector, and Entropy Bound}

\author{Sheng-Kai Huang}
\email{akai@fawstudio.com}
\affiliation{Independent Researcher}

\date{\today}

\maketitle

%=======================================================================
\section{S1. Quantum Eraser as Petz Recovery: Explicit Calculation}
\label{sec:eraser_calc}
%=======================================================================

We provide the complete calculation showing that the quantum eraser
phenomenon is precisely the Petz recovery map applied to the partial
trace channel.

\subsection{Setup}

Consider a two-path interferometer with which-path detector:

\emph{Initial state.}---The signal photon enters in a superposition of
paths $A$ and $B$:
\begin{equation}
\ket{\phi}_S = \frac{1}{\sqrt{2}}(\ket{A} + \ket{B}).
\end{equation}

\emph{Entangling interaction.}---A which-path detector entangles with the
signal:
\begin{equation}
\ket{\psi}_{SE} = \frac{1}{\sqrt{2}}(\ket{A}\ket{d_A}
  + \ket{B}\ket{d_B}),
\label{eq:entangled}
\end{equation}
where $\braket{d_A}{d_B} = 0$ (perfect which-path detection).

\emph{Which-path channel.}---The effective noise channel on the signal
is:
\begin{equation}
\chan(\rho_S) = \Tr_E\!\big[U(\rho_S \otimes \ket{d_0}\bra{d_0})U^\dagger
\big] = \sum_i \ket{i}\bra{i}\,\rho_S\,\ket{i}\bra{i},
\end{equation}
where $i \in \{A,B\}$. This is complete dephasing in the path basis.

\subsection{Computing the Petz Recovery}

We set the reference state on the full system $SE$:
\begin{equation}
\sigma_{SE} = \ket{\psi}\bra{\psi}_{SE}.
\end{equation}

The channel is $\chan = \Tr_E: \mathcal{B}(\hilb_S \otimes \hilb_E)
\to \mathcal{B}(\hilb_S)$. We compute:

\emph{Step 1: $\chan(\sigma_{SE})$.}
\begin{align}
\chan(\sigma_{SE}) &= \Tr_E\!\big[\ket{\psi}\bra{\psi}_{SE}\big]
\notag\\
&= \tfrac{1}{2}\big(\ket{A}\bra{A} + \ket{B}\bra{B}\big)
= \tfrac{1}{2}I_S.
\end{align}

\emph{Step 2: $\chan(\sigma_{SE})^{-1/2}$.}
\begin{equation}
\big(\tfrac{1}{2}I_S\big)^{-1/2} = \sqrt{2}\,I_S.
\end{equation}

\emph{Step 3: The adjoint $\chan^\dagger$.}
For $\chan = \Tr_E$, the Hilbert--Schmidt adjoint is defined by
$\Tr[\chan^\dagger(A) \cdot B] = \Tr[A \cdot \chan(B)]$ for all $A,B$.
One verifies $\Tr[(A \otimes I_E) \cdot B] = \Tr[A \cdot \Tr_E(B)]$,
so:
\begin{equation}
\Tr_E^\dagger(\rho_S) = \rho_S \otimes I_E.
\end{equation}

\emph{Step 4: Assembling the Petz map.}

For input $\rho_S$:
\begin{align}
\Petz_{\sigma,\Tr_E}(\rho_S)
&= \sigma_{SE}^{1/2}\;\Tr_E^\dagger\!\big(
  \sqrt{2}I_S\;\rho_S\;\sqrt{2}I_S
\big)\;\sigma_{SE}^{1/2} \notag\\
&= \ket{\psi}\bra{\psi}\;\big(2\rho_S \otimes
  I_E\big)\;\ket{\psi}\bra{\psi}.
\label{eq:petz_eraser}
\end{align}

Note: $\sigma_{SE}^{1/2} = \ket{\psi}\bra{\psi}$ since $\sigma_{SE}$
is a rank-1 projector.

\emph{Step 5: Evaluating for $\rho_S = \frac{1}{2}I_S$.}

Using orthogonality $\braket{d_A}{d_B} = 0$:
\begin{align}
&\bra{\psi}\big(2 \cdot \tfrac{1}{2}I_S \otimes I_E\big)\ket{\psi}
= \bra{\psi}(I_S \otimes I_E)\ket{\psi} = 1.
\end{align}

Therefore, the Petz map is trace-preserving and gives exact recovery:
\begin{equation}
\Petz_{\sigma,\Tr_E}(\rho_S) = \ket{\psi}\bra{\psi}_{SE}.
\end{equation}

This is exact recovery of the full entangled state, as expected for a
pure-state reference with the partial trace channel. The quantum eraser
measurement on $E$ (projecting onto
$\ket{\pm} = (\ket{d_A} \pm \ket{d_B})/\sqrt{2}$) then yields:
\begin{align}
\rho_S^{(+)} &= \tfrac{1}{2}(\ket{A}+\ket{B})(\bra{A}+\bra{B})
= \ket{\phi^+}\bra{\phi^+}, \\
\rho_S^{(-)} &= \tfrac{1}{2}(\ket{A}-\ket{B})(\bra{A}-\bra{B})
= \ket{\phi^-}\bra{\phi^-},
\end{align}
each with probability $1/2$. Interference is restored. \hfill$\square$

%=======================================================================
\section{S2. Self-Contained Proof of Theorem 1 (Uniqueness)}
\label{sec:proof_theorem1}
%=======================================================================

We provide a self-contained proof following
Parzygnat--Buscemi~\cite{Parzygnat2023}, emphasizing the key steps.

\subsection{The Classical Case}

For commutative $C^*$-algebras (i.e., classical probability), let
$p = (p_1, \ldots, p_n)$ be a prior and $N = (N_{j|i})$ a stochastic
matrix. Classical Bayes' rule gives:
\begin{equation}
R(N,p)_{i|j} = \frac{p_i\, N_{j|i}}{\sum_k p_k\, N_{j|k}}.
\label{eq:bayes}
\end{equation}

This can be written in the ``Petz form'':
\begin{equation}
R(N,p)(q)_i = p_i^{1/2}\;\sum_j \frac{N_{j|i}}{(Np)_j^{1/2}}
\cdot q_j \cdot \frac{1}{(Np)_j^{1/2}} \;\cdot p_i^{1/2}.
\end{equation}

In matrix notation, defining $D_p = \mathrm{diag}(p_1,\ldots,p_n)$:
\begin{equation}
R(N,p) = D_p^{1/2}\, N^T\, D_{Np}^{-1/2}\;(\cdot)\; D_{Np}^{-1/2}\,
N\, D_p^{1/2}.
\end{equation}

This is precisely the Petz map~(2) of the main text restricted to
diagonal matrices. The uniqueness of Bayes' rule in the classical
case is well-known.

\subsection{Extension to the Quantum Case}

The key insight is that axioms (ii) (functoriality) and (iv) (classical
limit) \emph{together} uniquely determine the noncommutative extension.

\emph{Step 1: Stinespring dilation.}---Any CPTP map $\chan$ can be
written as $\chan(\rho) = \Tr_E[V \rho V^\dagger]$ for some isometry
$V: \hilb_A \to \hilb_B \otimes \hilb_E$. The isometry $V$ is unique up
to unitary on $\hilb_E$.

\emph{Step 2: Factoring through the isometry.}---By functoriality
axiom (ii), $R(\chan, \sigma) = R(\Tr_E, V\sigma V^\dagger) \circ
R(\mathcal{V}, \sigma)$, where $\mathcal{V}(\rho) = V\rho V^\dagger$.
For the isometric channel $\mathcal{V}$, axioms (i) and (iii) imply
$R(\mathcal{V}, \sigma) = \mathcal{V}^{-1}$ (the adjoint
$V^\dagger(\cdot)V$). This is because an isometry has a unique left
inverse that is CPTP on the range.

\emph{Step 3: Reducing to the partial trace.}---It remains to determine
$R(\Tr_E, \omega)$ for $\omega = V\sigma V^\dagger$. The partial trace
is a product of marginalizations over each factor of $\hilb_E$. By
functoriality, $R(\Tr_E, \omega)$ factorizes accordingly.

\emph{Step 4: Classical limit forces the structure.}---For each factor
in the decomposition, axiom (iv) determines the map on the diagonal
(classical) part. The off-diagonal (quantum coherence) part is then
fixed by the requirement that $R$ be CP (complete positivity). The
unique CP extension of the classical Bayesian inverse is the Petz
map---this follows from the structure theory of CP maps on matrix
algebras.

The detailed argument uses the Choi--Jamio\l{}kowski isomorphism:
the Choi matrix of $R(\Tr_E, \omega)$ is uniquely determined by
its diagonal blocks (from axiom (iv)) and the requirement of positive
semidefiniteness of the Choi matrix (from axiom (i)). The unique
positive semidefinite completion is the Choi matrix of
$\Petz_{\omega, \Tr_E}$. \hfill$\square$

\subsection{What Happens When Each Axiom Is Dropped}

\begin{itemize}
\item \emph{Without (i) [CPTP]:} The map may not be physically
  realizable (not completely positive).
\item \emph{Without (ii) [functoriality]:} Multiple candidates exist,
  e.g., the ``rotated Petz map''~\cite{Wilde2015}
  $\tilde{R}(\rho) = \sigma^{1/2} U\, \chan^\dagger(\cdots)\,
  U^\dagger \sigma^{1/2}$ for arbitrary unitary $U$.
\item \emph{Without (iii) [normalization]:} One could rescale by an
  arbitrary positive function of $(\sigma, \chan)$.
\item \emph{Without (iv) [classical limit]:} The ``transpose channel''
  and other candidates remain viable.
\end{itemize}

%=======================================================================
\section{S3. Detailed Proof of Theorem 2 (QEC Optimality)}
\label{sec:proof_theorem2}
%=======================================================================

\subsection{Exact Recovery (Theorem 2a)}

Let $\code$ be a $[\![n,k,d]\!]$ stabilizer code with projector $P$,
and let $\chan(\rho) = \sum_a E_a \rho E_a^\dagger$ with Kraus
operators $\{E_a\}$ satisfying the Knill--Laflamme conditions:
\begin{equation}
P E_a^\dagger E_b P = c_{ab}\, P
\label{eq:KL}
\end{equation}
for a Hermitian matrix $C = (c_{ab})$.

Set $\sigma_\code = P/\dim\code$. Then for any $\rho$ supported on
$\code$:

\emph{Step 1:} $\chan(\sigma_\code) = \sum_a E_a (P/\dim\code)
E_a^\dagger$.

\emph{Step 2:} By the KL conditions, for $\rho$ supported on
$\code$:
\begin{equation}
D(\rho\|\sigma_\code) = D(\chan(\rho)\|\chan(\sigma_\code)).
\end{equation}
This is because the KL conditions ensure that the noise channel $\chan$
acts as an isometry on the code space (up to a decoupled environment),
preserving all relative entropies.

\emph{Step 3:} By Petz's sufficiency theorem~\cite{Petz1988}, equality
in the DPI implies $\Petz_{\sigma_\code,\chan} \circ \chan = \id$ on the
support of $\sigma_\code$, which is the code space $\code$. \hfill$\square$

\subsection{Near-Optimal Recovery (Theorem 2b)}

By Junge~et~al.~\cite{Junge2018}, for any $\sigma$ and $\chan$, there
exists a \emph{rotated} Petz map $\tilde{\Petz}$ such that:
\begin{equation}
F(\rho,\, \tilde{\Petz} \circ \chan(\rho))^2 \geq
\exp(-\Delta D),
\end{equation}
where $\Delta D = D(\rho\|\sigma) - D(\chan(\rho)\|\chan(\sigma))$.
This is the bound used in the main text (Eq.~(6)); it applies to the
\emph{rotated} Petz map, not the standard one.

The \emph{standard} Petz map $\Petz_{\sigma,\chan}$ satisfies the
weaker Barnum--Knill bound~\cite{Barnum2002}:
\begin{equation}
F_e(\code,\, \Petz_{\sigma_\code,\chan} \circ \chan) \geq
\frac{1}{2} F_e(\code,\, R_{\mathrm{opt}} \circ \chan)^2,
\end{equation}
where $F_e$ is the entanglement fidelity and $R_{\mathrm{opt}}$ is the
globally optimal decoder. Since $F_{e,\mathrm{opt}}^2 \geq
\exp(-\Delta D)$ by the Junge et~al.\ bound, this gives:
\begin{equation}
F_e(\code,\, \Petz_{\sigma_\code,\chan} \circ \chan) \geq
\tfrac{1}{2}\exp(-\Delta D).
\end{equation}
Note the factor of $1/2$: the standard Petz map is within a
constant factor of optimal, but does not saturate the exponential
bound in general.  For the saturation theorem
(Theorem~4 of the main text, $\sigma = I/d$), the bound
$F^2 \geq \exp(-\Delta D)$ for the standard Petz map is proved
directly via the AM--GM inequality (Sec.~S17), without invoking the
Junge et~al.\ machinery.

\subsection{Worked Example: $[\![5,1,3]\!]$ Code}

For the five-qubit code under depolarizing noise $\chan_p(\rho) =
(1-p)\rho + (p/3)(X\rho X + Y\rho Y + Z\rho Z)$ per qubit, the
Knill--Laflamme conditions are satisfied exactly for weight-$\leq 1$
errors. At physical error rate $p = 0.01$:

\begin{align}
F_{\mathrm{Petz}}^2 &\geq \exp(-\Delta D) \approx 1 - 5p^2 + O(p^3)
\approx 0.9995, \\
F_{\mathrm{opt}}^2 &\approx 0.9997.
\end{align}

The Petz decoder achieves 99.95\% fidelity vs.\ the optimal 99.97\%
---a gap of only $0.02\%$.

%=======================================================================
\section{S4. Detailed Proof of Proposition 3 (Master Inequality Chain)}
\label{sec:proof_theorem3}
%=======================================================================

We prove each link in the chain
$-\log F^2 \leq I(A;E|B) \leq \Sigma \leq \Delta D$ separately.

\subsection{Link 1: $-\log F^2 \leq \Delta D$ (Fidelity Bound)}

This is the Junge--Renner--Sutter--Wilde--Winter
theorem~\cite{Junge2018}. For the Petz recovery map:
\begin{equation}
F(\rho, \Petz_{\sigma,\chan}(\chan(\rho)))^2
\geq \exp\!\big(-(D(\rho\|\sigma) - D(\chan(\rho)\|\chan(\sigma)))\big).
\end{equation}

Taking logarithms: $-\log F^2 \leq \Delta D$.

\subsection{Link 2: $I(A;E|B) \geq -\log F^2$ (CMI Bound)}

Consider the Stinespring dilation of $\chan$:
$\chan(\rho_A) = \Tr_E[V \rho_A V^\dagger]$, producing the state
$\omega_{BE}$ from $\rho_A$. The conditional mutual information is:
\begin{equation}
I(A;E|B)_\omega = S(\omega_{AB}) + S(\omega_{BE}) - S(\omega_B) -
S(\omega_{ABE}).
\end{equation}

Fawzi--Renner~\cite{Fawzi2015} proved that for any tripartite state
$\omega_{ABE}$:
\begin{equation}
I(A;E|B)_\omega \geq -\log F(\omega_{AB},\,
\Petz_{B \to AB}(\omega_B))^2,
\end{equation}
where $\Petz_{B \to AB}$ is the Petz recovery map for the partial trace
$\Tr_A$. Combined with Link 1, this gives:
\begin{equation}
I(A;E|B) \geq -\log F^2.
\end{equation}

\subsection{Link 3: $\Sigma \geq I(A;E|B)$ (Entropy Production Bound)}

By the Kolchinsky--Wolpert decomposition~\cite{Kolchinsky2017}, the
entropy production of a process $\pi$ with initial distribution $p_0$
and steady state $q_0$ decomposes as:
\begin{align}
\Sigma(p_0) &= \Sigma_{\mathrm{Landauer}} + \Sigma_{\mathrm{mismatch}}
+ \Sigma_{\mathrm{res}} \\
&\geq \Sigma_{\mathrm{Landauer}} = D(p_0\|q_0) - D(\pi p_0\|\pi q_0).
\end{align}

The Landauer component $\Sigma_{\mathrm{Landauer}}$ captures the
information-theoretic cost.  For ideal (infinite, memoryless) thermal
baths, the total entropy production $\Sigma$ equals the relative entropy
decrease $\Delta D$; for finite-size baths,
$\Sigma_{\mathrm{Landauer}} \leq \Delta D$ and the mismatch and
residual terms provide finite-size corrections.
Since $\Delta D \geq I(A;E|B)$ (from Links 1--2) and
$\Sigma \geq \Sigma_{\mathrm{Landauer}} \geq 0$:
\begin{equation}
\Sigma \geq I(A;E|B).
\end{equation}
For ideal baths this follows from $\Sigma = \Delta D \geq I(A;E|B)$.
For finite baths, the chain $I(A;E|B) \leq \Sigma \leq \Delta D$
requires that the accessible entropy production exceeds the conditional
mutual information, which holds whenever the environment thermalizes
independently of the system's initial state (Markovian bath
assumption).

\subsection{Link 4: $\Delta D \geq \Sigma$ (Second Law)}

The relative entropy decrease $\Delta D = D(\rho\|\sigma) -
D(\chan(\rho)\|\chan(\sigma))$ is the total information lost under the
channel. By the Reeb--Wolf improved Landauer
principle~\cite{ReebWolf2014}, the accessible entropy production
$\Sigma$ satisfies:
\begin{equation}
\Delta D \geq \Sigma \geq 0.
\end{equation}
The gap $\Delta D - \Sigma$ accounts for system-bath correlations
established during the process.  For ideal (Markovian) baths,
$\Sigma = \Delta D$ and the chain collapses to
$-\log F^2 \leq I(A;E|B) \leq \Delta D$.

\subsection{Tightness Conditions}

\begin{center}
\begin{tabular}{lll}
\toprule
\textbf{Condition} & \textbf{Physical meaning} & \textbf{Domain} \\
\midrule
$F^2 = 1$ & Exact recovery & QEC \\
$I(A;E|B) = 0$ & Markov condition & No leakage \\
$\Sigma = 0$ & Reversible process & Thermo. \\
$\Delta D = 0$ & DPI saturated & Sufficiency \\
\bottomrule
\end{tabular}
\end{center}

\subsection{Gibbs State Specialization}

When $\sigma = \sigma_\beta = e^{-\beta H}/Z$:
\begin{align}
D(\rho\|\sigma_\beta) &= \Tr[\rho \log \rho] - \Tr[\rho \log
\sigma_\beta] \notag\\
&= -S(\rho) + \beta\langle H\rangle_\rho + \log Z \notag\\
&= \beta\big(\langle H\rangle_\rho - T S(\rho) - \mathcal{F}_{\mathrm{eq}}\big)
\notag\\
&= \beta\big(\mathcal{F}(\rho) - \mathcal{F}_{\mathrm{eq}}\big),
\end{align}
where $\mathcal{F}(\rho) = \langle H\rangle_\rho - T S(\rho)$ is the
non-equilibrium free energy and $\mathcal{F}_{\mathrm{eq}} = -\kB T \log Z$.
(We use $\mathcal{F}$ for free energy to avoid confusion with the
fidelity~$F$.)

Therefore:
\begin{align}
\Delta D &= D(\rho\|\sigma_\beta) - D(\chan(\rho)\|\chan(\sigma_\beta))
\notag\\
&= \beta\big[\mathcal{F}(\rho) - \mathcal{F}(\chan(\rho))\big]
  - \beta\big[\mathcal{F}_{\mathrm{eq}} - \mathcal{F}'_{\mathrm{eq}}\big] \notag\\
&= \beta\, W_{\mathrm{diss}}.
\end{align}

The master inequality becomes:
\begin{equation}
\boxed{
F^2 \geq \exp(-\beta\, W_{\mathrm{diss}})
= \exp(-\Delta S_{\mathrm{irr}}/\kB).
}
\end{equation}

This is a \emph{recovery-theoretic second law}: the ability to recover
a quantum state from a thermodynamic process is exponentially bounded by
the dissipated work.

%=======================================================================
\section{S5. Proof of Superposition Necessity}
\label{sec:superposition}
%=======================================================================

\begin{theorem}[Superposition from Recoverability]
If a system is in a definite-path state $\rho_S = \ket{i}\bra{i}$, then
for any recovery procedure (Petz map + measurement on the detector $E$),
the recovered signal state has no off-diagonal coherence in the path
basis. Consequently, recovery of interference implies prior
superposition.
\end{theorem}

\emph{Proof.}---

\emph{Step 1: No entanglement from definite paths.}
If $\rho_S = \ket{i_0}\bra{i_0}$ (definite path), the entangling
interaction produces:
\begin{equation}
\rho_{SE} = \ket{i_0}\bra{i_0} \otimes \ket{d_{i_0}}\bra{d_{i_0}}.
\end{equation}
This is a product state---no entanglement is created.

\emph{Step 2: Petz recovery on product states.}
The Petz recovery $\Petz_{\sigma_{SE}, \Tr_E}$ with
$\sigma_{SE} = \rho_{SE}$ (the product state) gives:
\begin{equation}
\Petz(\rho_S) = \ket{i_0}\bra{i_0} \otimes \ket{d_{i_0}}\bra{d_{i_0}},
\end{equation}
which is again a product state.

\emph{Step 3: Measurement on $E$ cannot create coherence.}
Any measurement $\{M_k\}$ on $E$ followed by post-selection gives:
\begin{align}
\rho_S^{(k)} &= \frac{\Tr_E[(I_S \otimes M_k)\,\rho_{SE}\,
  (I_S \otimes M_k^\dagger)]}{p_k} \notag\\
&= \frac{|\bra{d_{i_0}}M_k^\dagger\ket{d_{i_0}}|^2}{p_k}\,
  \ket{i_0}\bra{i_0} = \ket{i_0}\bra{i_0}.
\end{align}

No coherence. No interference. \hfill$\square$

\emph{Contrapositive:} If the eraser procedure produces
$\rho_S^{(k)}$ with $\bra{j}\rho_S^{(k)}\ket{l} \neq 0$ for $j \neq l$
(off-diagonal coherence = interference), then the initial state was
\emph{not} $\ket{i_0}\bra{i_0}$ for any $i_0$. It must have been a
superposition.

%=======================================================================
\section{S6. Quantum Crooks Fluctuation Theorem via Petz Map}
\label{sec:crooks}
%=======================================================================

Following Kwon--Kim~\cite{Kwon2019}, consider a thermodynamic protocol:

\emph{Initial state:} $\sigma_\beta^{(i)} = e^{-\beta H_i}/Z_i$.

\emph{Forward process:} unitary $U$ (changing $H_i$ to $H_f$), followed
by thermalization $\mathcal{T}$. Total: $\chan_F = \mathcal{T} \circ
\mathcal{U}$.

\emph{Reverse process:} Defined via the Petz map:
$\chan_R = \Petz_{\sigma_\beta^{(i)}, \chan_F}$.

By functoriality (Theorem 1, axiom (ii)):
\begin{equation}
\chan_R = \Petz_{\sigma_\beta^{(i)}, \mathcal{U}} \circ
\Petz_{U\sigma_\beta^{(i)}U^\dagger, \mathcal{T}}.
\end{equation}

Since $\Petz_{\sigma, \mathcal{U}} = U^\dagger(\cdot)U$ (the Petz
reverse of a unitary is its inverse):
\begin{equation}
\chan_R = \mathcal{U}^{-1} \circ
\Petz_{U\sigma_\beta^{(i)}U^\dagger, \mathcal{T}}.
\end{equation}

In the two-point measurement (TPM) scheme, the work distribution is:
\begin{equation}
p_F(w) = \sum_{n,m} \delta(w - E_m^f + E_n^i)\;
\frac{e^{-\beta E_n^i}}{Z_i}\;\abs{\bra{m_f}U\ket{n_i}}^2.
\end{equation}

For the reverse process defined by $\chan_R$:
\begin{equation}
p_R(-w) = \sum_{n,m} \delta(w - E_m^f + E_n^i)\;
\frac{e^{-\beta E_m^f}}{Z_f}\;\abs{\bra{n_i}U^\dagger\ket{m_f}}^2.
\end{equation}

Taking the ratio:
\begin{equation}
\frac{p_F(w)}{p_R(-w)} = \frac{Z_f}{Z_i}\, e^{\beta w}
= e^{\beta(w - \Delta F)},
\end{equation}
where $\Delta F = -\kB T \log(Z_f/Z_i)$. This is the Crooks
fluctuation theorem. \hfill$\square$

The Petz map naturally produces the correct time-reversed thermodynamic
process. The $e^{-\beta H/2}$ factors in the Petz map implement the
Boltzmann weighting that distinguishes thermodynamic time reversal from
mere unitary reversal.

%=======================================================================
\section{S7. Specialization Diagram}
\label{sec:diagram}
%=======================================================================

The four domains are connected by specialization of the pair
$(\sigma, \chan)$:

\begin{center}
\begin{tabular}{ccc}
(ii) Time Reversal & $\xrightarrow{\sigma = P/|\code|}$ & (iii) QEC \\
$\downarrow_{\chan = \Tr_E}$ & & $\downarrow_{|\code|=1}$ \\
(i) Quantum Eraser & $\xleftarrow{\code = \{\ket{\phi}\}}$ &
  (iii) single codeword \\
$\downarrow_{\sigma = \sigma_\beta}$ & & \\
(iv) Thermodynamics & &
\end{tabular}
\end{center}

\emph{QEC with single codeword = Quantum eraser.}---For a
one-dimensional code $\code = \mathrm{span}\{\ket{\phi}\}$, the
Knill--Laflamme conditions $P E_a^\dagger E_b P = c_{ab} P$ are
\emph{automatically satisfied} (every operator has a matrix element
in a one-dimensional space). Therefore, exact Petz recovery is always
possible for a single codeword. This is the quantum eraser: the ``code''
is the coherent superposition state, and the ``noise'' is decoherence.

\emph{Universal bound connects all.}---The bound
$F^2 \geq \exp(-\Delta D)$ specializes to:
\begin{itemize}
\item \emph{Quantum eraser:} $\Delta D = S(\rho_S^{\mathrm{decoh}}) -
  S(\rho_S^{\mathrm{init}})$ (entropy of decoherence).
  $\Delta D = 0$ for pure bipartite state $\Rightarrow$ exact recovery.
\item \emph{QEC:} $\Delta D = D(\rho\|\sigma_\code) -
  D(\chan(\rho)\|\chan(\sigma_\code))$ (information lost to noise).
  $\Delta D = 0$ when KL conditions hold.
\item \emph{Thermodynamics:} $\Delta D = \beta W_{\mathrm{diss}}$
  (dissipated work). $F^2 \geq e^{-\beta W_{\mathrm{diss}}}$.
\end{itemize}

%=======================================================================
\section{S8. Notation and Conventions}
\label{sec:notation}
%=======================================================================

\begin{center}
\begin{tabular}{ll}
\toprule
\textbf{Symbol} & \textbf{Meaning} \\
\midrule
$\hilb$ & Hilbert space \\
$\mathcal{B}(\hilb)$ & Bounded operators on $\hilb$ \\
$\mathcal{S}(\hilb)$ & Density matrices on $\hilb$ \\
$\chan$ & Quantum channel (CPTP map) \\
$\chan^\dagger$ & Hilbert--Schmidt adjoint \\
$\Petz_{\sigma,\chan}$ & Petz recovery map \\
$D(\rho\|\sigma)$ & Quantum relative entropy \\
$F(\rho,\sigma)$ & Uhlmann fidelity \\
$I(A;B|C)$ & Conditional mutual information \\
$\Sigma$ & Entropy production \\
$\sigma_\beta$ & Gibbs state $e^{-\beta H}/Z$ \\
$\kB$ & Boltzmann constant \\
$P$ & Code projector \\
$\code$ & Code space \\
\bottomrule
\end{tabular}
\end{center}

\emph{Relative entropy:} $D(\rho\|\sigma) = \Tr[\rho(\log\rho -
\log\sigma)]$, with the convention $D(\rho\|\sigma) = +\infty$ if
$\mathrm{supp}(\rho) \not\subseteq \mathrm{supp}(\sigma)$.

\emph{Fidelity:} $F(\rho,\sigma) = \norm{\sqrt{\rho}\sqrt{\sigma}}_1
= \Tr\sqrt{\sqrt{\rho}\sigma\sqrt{\rho}}$.

%=======================================================================
\section{S9. Post-Selection Analysis}
\label{sec:postselection}
%=======================================================================

We derive the post-selection improvement factor stated in the main text
and connect it to the Ising--QEC mapping.

\subsection{Derivation of the Improvement Factor}

Consider a stabilizer code $\code$ of distance $d$ under noise
$\chan$. Syndrome measurement yields outcome $s$ with probability
$p(s)$. For each syndrome, define the instance-level data-processing gap:
\begin{equation}
\Delta D(s) = D(\rho^{(s)}\|\sigma_\code)
  - D(\chan(\rho^{(s)})\|\chan(\sigma_\code)),
\end{equation}
where $\rho^{(s)}$ is the post-syndrome logical state. By the
universal recovery bound, $F(s)^2 \geq \exp(-\Delta D(s))$.

The overall logical failure probability is:
\begin{equation}
p_{\mathrm{fail}} = \sum_s p(s)\, \big(1 - F(s)^2\big)
\leq \sum_s p(s)\, \big(1 - e^{-\Delta D(s)}\big).
\end{equation}

Post-selecting: accept only syndromes in
$\mathcal{A}_\varepsilon = \{s : \Delta D(s) \leq
\Delta D_{\mathrm{th}}(\varepsilon)\}$,
where $\sum_{s \notin \mathcal{A}_\varepsilon} p(s) = \varepsilon$.
The post-selected failure probability is:
\begin{equation}
p_{\mathrm{fail}}^{(\varepsilon)}
= \frac{1}{1-\varepsilon}\sum_{s \in \mathcal{A}_\varepsilon} p(s)\,
  \big(1 - F(s)^2\big)
\leq \Delta D_{\mathrm{th}}(\varepsilon),
\end{equation}
where the last step uses $1 - e^{-x} \leq x$ for $x \geq 0$.

The improvement factor is:
\begin{equation}
R(\varepsilon) = \frac{p_{\mathrm{fail}}}{p_{\mathrm{fail}}^{(\varepsilon)}}
\approx \exp\!\big(\Delta D_{\mathrm{th}}(\varepsilon)\big),
\end{equation}
where the exponential form holds when the dominant contribution to
failures comes from the tail of the $\Delta D$ distribution.

\subsection{Distribution of $\Delta D(s)$}

For a distance-$d$ surface code under depolarizing noise at rate $p$,
errors of weight $w$ have probability $\sim p^w$ and produce
$\Delta D(s) \approx w \ln(1/p)$. The tail probability is:
\begin{equation}
\Pr[\Delta D > x] \sim \exp(-x/\lambda),
\quad \lambda \sim \ln(1/p).
\end{equation}
This exponential tail makes post-selection effective: a small fraction
of high-$\Delta D$ instances carries a disproportionate share of
failures.

\subsection{Connection to the Ising--QEC Mapping}

English~et~al.~\cite{EnglishEtAl2025} established an exact mapping
from surface-code failure rates to Ising-model partition functions:
\begin{equation}
p_{\mathrm{fail}}(s) = \frac{Z_{-}(s)}{Z_{+}(s) + Z_{-}(s)}
= \frac{1}{1 + e^{\Delta F(s)/T_{\mathrm{eff}}}},
\end{equation}
where $\Delta F(s)$ is the Ising domain-wall free energy. Identifying
$\Delta F(s)/T_{\mathrm{eff}} \approx \Delta D(s)$, the Petz bound
$F(s)^2 \geq e^{-\Delta D(s)}$ becomes a statement about Ising
free-energy gaps: post-selection on low-$\Delta D$ syndromes is
equivalent to conditioning on Ising configurations with large
free-energy gaps, where the correct logical value is
thermodynamically favored.

\subsection{The Scaling Prediction}

For a distance-$d$ code at error rate $p$ with rejection fraction
$\varepsilon$:
\begin{equation}
\ln R(\varepsilon, d) = \alpha(p)\, d + \beta(\varepsilon),
\end{equation}
where $\alpha(p) \approx \tfrac{1}{2}\ln(p_{\mathrm{th}}/p)$ (the
per-unit-distance improvement, decoder-independent) and
$\beta(\varepsilon) \approx \ln(1/(1-\varepsilon))$ (the rejection
overhead, code-distance-independent). The decoder-independence of
$\alpha$ follows from $\Delta D(s)$ depending only on $(\chan,
\sigma_\code)$, not on the recovery map.

This prediction is supported by: quadratic post-selection
gains~\cite{SmithBrownBartlett2024}, Ising-based
thresholds~\cite{EnglishWilliamsonBartlett2025}, scalable gains
$p_f \to p_f^b$ with $b \geq 2$~\cite{ChenEtAl2025}, and
below-threshold surface code operation~\cite{Acharya2025}.

%=======================================================================
\section{S10. Decoder Retrodiction Hierarchy}
\label{sec:decoder_hierarchy}
%=======================================================================

We develop the interpretation of the QEC decoder hierarchy as a
sequence of approximations to optimal Petz retrodiction.

\subsection{ML Decoding as Exact Retrodiction}

For a stabilizer code with syndrome $s$, maximum-likelihood decoding
selects:
\begin{equation}
\ell^* = \arg\max_\ell \; p(\ell \,|\, s)
= \arg\max_\ell \; p(s \,|\, \ell)\, p(\ell).
\end{equation}
This is classical Bayes' rule---the classical limit (axiom (iv) of
Theorem~1) of the Petz retrodiction functor. For the full quantum
problem, the Petz map $\Petz_{\sigma_\code,\chan}$ generalizes ML
by retrodicting the full quantum state including coherences.

\subsection{The Retrodiction Gap Hierarchy}

For decoder $\mathcal{D}$, define:
\begin{equation}
\delta_{\mathcal{D}} = D\!\left(\tilde{\Petz}_{\sigma_\code,\chan}
  \;\big\|\; \mathcal{D}\right),
\end{equation}
measuring the distinguishability from optimal retrodiction.
The established hierarchy~\cite{deMartiOlius2024} becomes:
\begin{equation}
0 = \delta_{\mathrm{ML}}
< \delta_{\mathrm{NN}}
< \delta_{\mathrm{CM}}
< \delta_{\mathrm{MWPM}}
< \delta_{\mathrm{UF}}.
\end{equation}

The bound from the main text gives:
\begin{equation}
F(\rho,\, \mathcal{D} \circ \chan(\rho))^2
\geq F(\rho,\, \tilde{\Petz} \circ \chan(\rho))^2
\cdot e^{-\delta_{\mathcal{D}}}.
\end{equation}

\subsection{The Computational Entropy Conjecture}

We conjecture $\delta_{\mathcal{D}} \leq c\,\sigma_{\mathcal{D}}$,
where $\sigma_{\mathcal{D}}$ is the thermodynamic entropy produced
during the decoder's computation. The motivation is Landauer's
principle: erasing correlations during decoding costs
$\kB T \ln 2$ per bit and simultaneously degrades retrodictive
information.

\emph{Heuristic support:}
\begin{itemize}
\item Union-Find: greedy local rules discard long-range correlations
  (high $\sigma$, large $\delta$).
\item MWPM: global matching but ignores error correlations
  (moderate $\sigma$ and $\delta$).
\item Neural networks (AlphaQubit~\cite{Bausch2024}): learn
  correlations from data, potentially lower $\delta$.
\item ML/tensor-network: full probabilistic inference,
  minimal information erasure ($\delta \approx 0$).
\end{itemize}

\subsection{Regime of Applicability}

Preliminary numerical experiments on $d = 3$ surface codes show that
decoder \emph{expressiveness} (the ability to represent complex
correlations) dominates over thermodynamic efficiency in the small-code
regime. A highly dissipative but expressive neural network outperforms
a thermodynamically efficient but inexpressive decoder. We expect the
conjecture to become relevant at $d \gtrsim 15$--$20$, where the
syndrome space grows exponentially and efficient information
compression becomes the bottleneck.

Even without the thermodynamic conjecture, the retrodiction-gap
framework provides independent value: $\delta_{\mathcal{D}}$ is a
decoder-independent metric computable from output statistics,
offering a unified lens for comparing decoders of entirely different
architectures.

\emph{Caveat on current evidence.}---The rank-order consistency
between the retrodiction hierarchy and known decoder performance
reported in the main text is a \emph{prediction}, not a verified
result. Computing $\delta_{\mathcal{D}}$ from published data requires
access to the full CPTP decoder map, which is not available for most
decoders. Direct numerical verification of
$\delta_{\mathrm{NN}} < \delta_{\mathrm{MWPM}} < \delta_{\mathrm{UF}}$
from decoder output statistics is an important open task.

%=======================================================================
\section{S11. Temporal Asymmetry and the Equivalence Chain}
\label{sec:equiv_chain}
%=======================================================================

We prove the equivalence chain stated in the main text and develop
the temporal asymmetry parameter $\tauasym$.

\subsection{Definition and Basic Properties}

For a quantum channel $\chan$, faithful reference state $\sigma$, and
input $\rho$, define:
\begin{equation}
\tauasym(\rho, \sigma, \chan) = 1 - F\!\big(\rho,\,
  \tilde{\Petz}_{\sigma,\chan}(\chan(\rho))\big),
\end{equation}
where $\tilde{\Petz}$ is the rotated Petz map of
Junge~et~al.~\cite{Junge2018}.

\emph{Basic properties:}
\begin{enumerate}
\item $0 \leq \tauasym \leq 1$ (from $0 \leq F \leq 1$).
\item $\tauasym = 0$ if and only if exact Petz recovery:
  $\tilde{\Petz}_{\sigma,\chan} \circ \chan(\rho) = \rho$.
\item $\tauasym$ is continuous in all arguments (by continuity of
  fidelity).
\item $\tauasym$ is monotone under composition: for composable
  $\chan_1, \chan_2$,
  $\tauasym(\rho, \sigma, \chan_2 \circ \chan_1) \geq
  \tauasym(\rho, \sigma, \chan_1)$, since additional processing
  cannot improve recovery.
\end{enumerate}

\subsection{The Equivalence Chain: Proof}

We prove that for a faithful $\sigma$ and channel $\chan$ with
Stinespring dilation $V: \hilb_A \to \hilb_B \otimes \hilb_E$, the
following are equivalent:

\noindent
(a) $\tauasym = 0$ (exact Petz recovery). \\
(b) $\Delta D = 0$ (DPI saturated). \\
(c) $I(A;E|B)_\omega = 0$ (quantum Markov chain). \\
(d) $\Sigma = 0$ (zero entropy production, for thermal $\sigma$). \\
(e) Detailed balance: $\Petz_{\sigma,\chan} = \chan$ (the channel
    is its own Petz reversal).

\emph{(a) $\Leftrightarrow$ (b):} This is the Petz sufficiency
theorem~\cite{Petz1988}. Exact recovery $F = 1$ holds if and only if
$D(\rho\|\sigma) = D(\chan(\rho)\|\chan(\sigma))$ for all $\rho$ in
the support of $\sigma$, i.e., $\Delta D = 0$.

\emph{(b) $\Leftrightarrow$ (c):} By Fawzi--Renner~\cite{Fawzi2015},
$I(A;E|B)_\omega = 0$ if and only if $\omega_{ABE}$ is a quantum
Markov chain $A - B - E$. By the Stinespring structure, this is
equivalent to perfect recoverability of $A$ from $B$, i.e.,
$\Delta D = 0$.

\emph{(b) $\Leftrightarrow$ (d) (for Gibbs $\sigma$):} When $\sigma =
e^{-\beta H}/Z$, the entropy production $\Sigma = \beta W_{\mathrm{diss}}
+ D_{\mathrm{mismatch}}$ by Kolchinsky--Wolpert~\cite{Kolchinsky2017}.
$\Sigma = 0$ requires both $W_{\mathrm{diss}} = 0$ and zero mismatch,
which for the Gibbs reference means $\Delta D = \beta W_{\mathrm{diss}}
= 0$. The converse follows from $\Sigma \leq \Delta D$ (second
law, Reeb--Wolf~\cite{ReebWolf2014}).

\emph{(d) $\Leftrightarrow$ (e):} By Kwon--Kim~\cite{Kwon2019}:
the Crooks ratio $P_F/P_R = \exp(\Sigma/\kB)$. Detailed balance
($\Petz_{\sigma,\chan} = \chan$, i.e., the channel equals its own
reversal) holds if and only if $P_F = P_R$ for all trajectories, i.e.,
$\Sigma = 0$.

This completes the chain. \hfill$\square$

\subsection{Bound: $\tauasym \leq 1 - \exp(-\Sigma/2)$}

From Proposition~3 of the main text (master inequality chain):
\begin{equation}
F^2 \geq \exp(-\Sigma).
\end{equation}
Since $F \geq 0$ and $\exp(-x/2) \geq 0$, taking the square root:
\begin{equation}
F \geq \exp(-\Sigma/2).
\end{equation}
Therefore:
\begin{equation}
\tauasym = 1 - F \leq 1 - \exp(-\Sigma/2).
\end{equation}

For small $\Sigma$: Taylor-expanding gives $\tauasym \lesssim \Sigma/2$.
For large $\Sigma$: $\tauasym \to 1$. The bound is tight when
$F^2 = \exp(-\Sigma)$ (saturation of the Junge bound).

\subsection{Connection to the Crooks Fluctuation Theorem}

The Crooks theorem~\cite{Crooks1999} states:
\begin{equation}
\frac{P_F(w)}{P_R(-w)} = \exp\!\big(\beta(w - \Delta F)\big).
\end{equation}
For a cyclic process ($\Delta F = 0$), averaging over trajectories:
$\langle P_F/P_R \rangle = \exp(\langle \beta W \rangle) =
\exp(\Sigma/\kB)$. When $\Sigma = 0$: $P_F = P_R$ pointwise.
Forward and reverse processes are equally probable for every
trajectory.

The temporal asymmetry parameter $\tauasym$ provides a
\emph{state-level} (rather than trajectory-level) measure of this
asymmetry. While the Crooks ratio is defined over stochastic
trajectories, $\tauasym$ is defined for a single input state $\rho$
and measures whether that state can be recovered after the channel.
The connection is: $\tauasym = 0$ implies $\Sigma = 0$ implies
$P_F = P_R$ (the process is its own time reversal).

This is the precise mathematical content of the physical insight:
in a zero-entropy-production system, there is no statistical
distinction between the forward and reverse time directions.

\subsection{The Retrodiction Landauer Principle}

The energy cost of retrodiction---inferring the pre-channel state
from the post-channel state---is bounded by the entropy produced
during the forward process.

\emph{Observation (Retrodiction Landauer principle).}---For a
channel $\chan$ with thermal reference state
$\sigma_\beta = e^{-\beta H}/Z$, we argue heuristically that the
minimum energy cost of any physical implementation of approximate
retrodiction satisfies:
\begin{equation}
E_{\mathrm{retro}} \geq \kB T \cdot \Sigma.
\label{eq:retro_landauer}
\end{equation}
The argument applies Landauer's principle~\cite{ReebWolf2014}
(energy dissipation $\geq \kB T \ln 2$ per bit erased) to the
task of recovering $\Delta D$ nats of distinguishability from
$\chan(\rho)$ to $\rho$, giving
$E_{\mathrm{retro}} \geq \kB T \cdot \Delta D$.
By the master chain, $\Delta D \geq \Sigma$, giving
Eq.~\eqref{eq:retro_landauer}.
This extends Landauer's principle from bit erasure to state
retrodiction; the extension is physically motivated but not
rigorously derived from first principles.

\emph{Physical consequences:}
\begin{itemize}
\item $\Sigma = 0$ (closed system): retrodiction is
  thermodynamically free. The past is as accessible as the future.
\item $\Sigma > 0$ (open system): retrodiction has a minimum energy
  cost proportional to the entropy produced. The past becomes
  energetically expensive to access.
\item The bound is tight when $\Delta D = \Sigma$ (all information
  loss is dissipative). \hfill$\square$
\end{itemize}

\subsection{The Quantum Eraser as the $\tauasym = 0$ Special Case}

For a pure bipartite state $\ket{\psi}_{SE}$ with $\chan = \Tr_E$:

\begin{itemize}
\item $\sigma = \ket{\psi}\bra{\psi}_{SE}$ (pure).
\item $D(\sigma\|\sigma) = 0$ and
  $D(\chan(\sigma)\|\chan(\sigma)) = 0$, so $\Delta D = 0$ trivially.
\item By the equivalence chain: $\Delta D = 0 \Rightarrow
  I(A;E|B) = 0 \Rightarrow \Sigma = 0 \Rightarrow \tauasym = 0$.
\item $F = 1$: exact recovery, as computed in Sec.~S1.
\end{itemize}

The quantum eraser operates in the $\tauasym = 0$ regime. The idler
measurement ``works'' not because of backward causation but because
the signal-idler pair forms a closed system with zero entropy
production. In such a system, the Crooks ratio is exactly 1, and
there is no statistical distinction between measuring the signal
first or the idler first.

The ``need to wait'' for the idler measurement arises only in the
\emph{open-system description}: when we trace out the idler and
describe only the signal, we have $\tauasym > 0$ for the signal
alone. The full bipartite system has $\tauasym = 0$, but the
subsystem has $\tauasym > 0$. Recovering $\tauasym = 0$ requires
accessing the idler---this is the operational meaning of ``waiting
for the future measurement.''

%=======================================================================
\section{S12. Numerical Verification: Design Rationale and Results}
\label{sec:numerical}
%=======================================================================

In this section we describe numerical experiments designed to test the
\emph{new predictions} of our framework (Predictions~1 and~2 of the
main text), and explain why these experiments go beyond verifying
known mathematics.

\subsection{Scope of Numerical Tests}

The individual inequalities in the master chain~(8) of the main text
are rigorously proven~\cite{Junge2018,Fawzi2015,ReebWolf2014};
computing them for specific channels confirms implementation
correctness but does not test any physical prediction.
The numerical experiments below instead test consequences that follow
from the \emph{uniqueness} of the Petz map as a
retrodiction functor (Theorem~1) and from the \emph{structure} of
the master inequality chain when applied to error correction.

\subsection{Experiment~1: Decoder Retrodiction Hierarchy}

\emph{Known:} Different QEC decoders (ML, MWPM, Union-Find, etc.)
have different logical error rates~\cite{deMartiOlius2024}. This
performance ordering is well-established empirically.

\emph{New prediction:} The retrodiction gap
$\delta_{\mathcal{D}} = D(\tilde{\Petz}_{\sigma,\chan}\|\mathcal{D})$
should \emph{quantitatively rank-order} decoder performance.
Specifically, decoders with smaller $\delta_{\mathcal{D}}$ must have
higher recovery fidelity.

\emph{Why this is not obvious:} If the Petz map were merely
\emph{one of many} good recovery maps, then
$\delta_{\mathcal{D}}$ (the divergence from the Petz map) would have
no reason to predict decoder quality --- a decoder could differ
greatly from the Petz map yet still perform well by approximating
some other good recovery map. The prediction that
$\delta_{\mathcal{D}}$ rank-orders performance follows specifically
from the \emph{uniqueness} of the Petz map (Theorem~1): since
it is the \emph{only} Bayesian-consistent retrodiction functor, any
sufficiently good decoder must converge toward it.

\emph{Numerical protocol:}
\begin{enumerate}
\item Fix a surface code of distance $d$ under depolarizing noise
  at rate $p$.
\item Implement four decoders: maximum-likelihood (ML, via exhaustive
  enumeration for small $d$), minimum-weight perfect matching (MWPM),
  Union-Find (UF), and a simple lookup-table decoder.
\item For each decoder $\mathcal{D}$, compute the effective
  channel $\mathcal{D} \circ \chan$ by Monte Carlo sampling over
  error configurations.
\item Compute
  $\delta_{\mathcal{D}} = D(\tilde{\Petz}_{\sigma,\chan}\|\mathcal{D})$
  via the channel divergence between the Petz decoder and
  $\mathcal{D}$.
\item Compare the ordering $\delta_{\mathrm{ML}} < \delta_{\mathrm{MWPM}}
  < \delta_{\mathrm{UF}} < \delta_{\mathrm{lookup}}$ against the
  known performance ordering.
\end{enumerate}

The prediction fails if a decoder with
\emph{larger} $\delta_{\mathcal{D}}$ consistently outperforms one
with \emph{smaller} $\delta_{\mathcal{D}}$, which would indicate that
the Petz map's uniqueness as a retrodiction functor does not
constrain the structure of practical decoders.

\subsection{Experiment~2: Post-Selection Thermodynamic Filtering}

\emph{Known:} Post-selection on syndrome outcomes can improve
logical error rates in QEC~\cite{SmithBrownBartlett2024,ChenEtAl2025}.
The improvement depends on the rejection fraction $\varepsilon$ and
the code distance~$d$.

\emph{New prediction:} The improvement factor obeys
$\ln R(\varepsilon, d) = \alpha(p)\, d + \beta(\varepsilon)$,
where:
\begin{itemize}
\item $\alpha(p)$ depends on the noise rate $p$ but is
  \emph{independent of the decoder} $\mathcal{D}$;
\item $\beta(\varepsilon)$ depends on the rejection fraction but is
  \emph{independent of the code distance} $d$.
\end{itemize}

\emph{Why $\alpha$ is decoder-independent:} The data-processing gap
$\Delta D(s)$ for syndrome $s$ is determined by the noise channel
$\chan$ and the code state $\sigma_\code$ alone --- it measures how
much distinguishability is lost under the noise, which is a property
of the noise and the code, not of the decoder. The per-distance
scaling $\alpha(p)$ inherits this decoder-independence because it is
set by the distribution of $\Delta D(s)$ values, which is a property
of the channel.

\emph{Numerical protocol:}
\begin{enumerate}
\item Fix surface codes of distances $d = 3, 5, 7$ under
  depolarizing noise at rates $p \in \{0.01, 0.03, 0.05\}$.
\item For each $(d, p)$, sample $10^5$--$10^6$ syndrome outcomes and
  compute $\Delta D(s)$ for each.
\item Apply multiple decoders (MWPM, UF, lookup) to each syndrome.
\item For each decoder and rejection fraction $\varepsilon \in
  [0, 0.3]$, compute the post-selected logical error rate
  $p_{\mathrm{fail}}^{(\varepsilon)}$ and the improvement factor
  $R(\varepsilon, d)$.
\item Fit $\ln R(\varepsilon, d) = \alpha(p)\, d +
  \beta(\varepsilon)$ and verify:
  \begin{itemize}
  \item[(a)] The factorization holds (linear in $d$ with
    $d$-independent intercept);
  \item[(b)] $\alpha(p)$ is the same across different decoders.
  \end{itemize}
\end{enumerate}

A decoder-dependent $\alpha$ --- e.g., MWPM and Union-Find
giving different slopes at the same noise rate --- would indicate
that the thermodynamic filtering picture is incomplete.

\subsection{Relationship to the Quantum Eraser Calculation}

The explicit quantum eraser calculation (Sec.~S1) serves a different
purpose from the numerical experiments above. The eraser calculation
demonstrates that the Petz map formalism correctly reproduces the
known physics of quantum erasure ($\tauasym = 0$ for a closed
bipartite system). This is a \emph{consistency check} of the formalism,
not a test of a new prediction. The numerical experiments in this
section, by contrast, test consequences of the framework that are
\emph{not known a priori} and could in principle be falsified.

%=======================================================================
\section{S13. Experimental Proposals and Existing Evidence}
\label{sec:experiments}
%=======================================================================

The master inequality chain~(8) of the main text and the predictions
derived from it can be tested against both existing published data and
purpose-built experiments. In this section we describe reanalysis
strategies for published results (Sec.~S13.1), propose two new
experimental protocols (Secs.~S13.2--S13.3), and specify the
simulation methodology used in this work (Sec.~S13.4).

\subsection{S13.1: Reanalysis of Published Data}

Several recent experimental results can be reinterpreted through the
lens of the Petz retrodiction framework, providing immediate
consistency checks without new hardware.

\emph{Decoder hierarchy.}---Bausch et~al.~\cite{Bausch2024}
(AlphaQubit) demonstrated that a recurrent neural-network decoder
outperforms correlated minimum-weight perfect matching (MWPM), which
in turn outperforms standard MWPM, which outperforms Union-Find (UF).
Our framework (Sec.~S10) predicts that this performance ordering is
equivalent to the retrodiction divergence ordering:
$\delta_{\mathrm{NN}} < \delta_{\mathrm{cMWPM}}
< \delta_{\mathrm{MWPM}} < \delta_{\mathrm{UF}}$.
A concrete test is to compute $\delta_{\mathcal{D}}$ for each decoder
from the published logical error rates and noise model parameters
reported in Ref.~\cite{Bausch2024}. If the CPTP map of the decoder
is not fully specified, one can use an operational proxy:
$\delta_{\mathcal{D}}^{\mathrm{op}} = p_L(\mathcal{D}) -
p_L(\mathrm{ML})$, normalized by the code distance, where
$p_L(\mathrm{ML})$ is the maximum-likelihood logical error rate.

\emph{Post-selection.}---English et~al.~\cite{EnglishWilliamsonBartlett2025}
and Chen et~al.~\cite{ChenEtAl2025} demonstrated that post-selecting
on syndrome outcomes improves the effective QEC threshold. Our framework
predicts that the improvement factor satisfies
$\ln R = \alpha(p)\cdot d + \beta(\varepsilon)$ with the
per-distance coefficient $\alpha$ being decoder-independent
(Sec.~S9). Reanalyzing their published data to extract $\alpha$
for different decoders at the same noise rate would directly test
this prediction. Preliminary inspection of the data in
Ref.~\cite{ChenEtAl2025} shows scaling consistent with
decoder-independent $\alpha$, but a systematic extraction has not
been performed.

\emph{Petz recovery experiments.}---Singh et~al.~\cite{Singh2025}
implemented the Petz recovery map on a nuclear magnetic resonance
(NMR) platform, and Png and Scarani~\cite{Png2025} realized it on a
trapped-ion quantum processor. Both experiments measured aspects of
recovery fidelity under controlled noise channels.
Our framework predicts $F \geq \exp(-\Sigma/2)$, where $\Sigma$ is
the entropy production of the noise process. Two important caveats
apply: (i)~neither experiment directly measured the Petz recovery
fidelity $F$ in the precise sense of Eq.~(7) of the main text
(i.e., fidelity between the input state and the output of the
rotated Petz map); rather, they measured circuit-level recovery
fidelities that serve as proxies; and (ii)~neither experiment
simultaneously measured $\Sigma$. One can compute $\Sigma$ from
the reported channel parameters (e.g., amplitude damping with
measured decay rate $\gamma$, giving
$\Sigma(\gamma) = \gamma \ln[(1+\gamma)/\gamma]$ for a thermal bath)
and check consistency, but this constitutes an indirect test based on
the theoretical channel model, not a direct simultaneous measurement
of $F$ and $\Sigma$ on the same experimental run. A simultaneous
measurement of both quantities on the same quantum channel is an
important open experimental direction (see Sec.~S13.2).

\subsection{S13.2: Proposed Experiment---Simultaneous $F$ and $\Sigma$ Measurement}

To our knowledge, no existing experiment has simultaneously measured
the Petz recovery fidelity $F$ and the entropy production $\Sigma$ for
the same quantum channel in the same experimental run.
This represents a significant gap in the experimental evidence base:
all existing tests of the bound $F \geq \exp(-\Sigma/2)$ rely either on
theoretical calculations from channel parameters (Sec.~S15.1) or on
separate measurements of $F$ and $\Sigma$ in different experimental
contexts. A simultaneous measurement would provide the most direct
test of both the bound and the equivalence $\tauasym = 0
\Leftrightarrow \Sigma = 0$.

\emph{System.}---A superconducting transmon qubit coupled to a tunable
electromagnetic environment (e.g., a lossy resonator with adjustable
coupling~$\kappa$). The effective noise channel is amplitude damping
$\chan_\gamma$ with damping parameter
$\gamma = 1 - e^{-\kappa t}$ controlled by the coupling time~$t$.

\emph{Protocol:}
\begin{enumerate}
\item \emph{Preparation}: Initialize the qubit in a known state
  $\rho$ (e.g., $\ket{+} = (\ket{0}+\ket{1})/\sqrt{2}$) via
  standard gate operations.
\item \emph{Noise application}: Apply the amplitude damping channel
  $\chan_\gamma$ by coupling to the lossy resonator for time~$t$.
  Vary $\gamma$ from $0$ to $1$ by sweeping $t$.
\item \emph{Measurement~1 (entropy production)}: Perform quantum state
  tomography on $\chan_\gamma(\rho)$ to reconstruct the output state.
  Compute $\Sigma$ from the channel parameters:
  \begin{equation}
  \Sigma(\gamma) = D\!\big(\chan_\gamma(\rho)\big\|\chan_\gamma(\sigma_\beta)\big)
  - D(\rho\|\sigma_\beta) + \beta\, Q,
  \end{equation}
  where $Q = \Tr[H(\rho - \chan_\gamma(\rho))]$ is the heat
  dissipated and $\sigma_\beta$ is the thermal state of the qubit.
  For a zero-temperature bath, $\Sigma(\gamma) = (1-\gamma)\ln(1-\gamma)
  + \gamma$ (in nats) when $\rho = \ket{1}\bra{1}$.
\item \emph{Measurement~2 (recovery fidelity)}: Implement the Petz
  recovery map $\Petz_{\sigma,\chan_\gamma}$ as a quantum circuit.
  For amplitude damping, the Petz map has a known circuit
  decomposition~\cite{Barnum2002} involving an ancilla qubit
  and controlled rotations. Apply $\Petz_{\sigma,\chan_\gamma}$ to
  $\chan_\gamma(\rho)$ and perform state tomography on the output.
  Compute $F = F(\rho, \Petz_{\sigma,\chan_\gamma}(\chan_\gamma(\rho)))$.
\item \emph{Sweep}: Repeat steps 1--4 for $\gamma \in \{0, 0.1, 0.2,
  \ldots, 1.0\}$.
\end{enumerate}

\emph{Predictions:}
\begin{itemize}
\item $F(\gamma) \geq \exp(-\Sigma(\gamma)/2)$ for all $\gamma$.
\item At $\gamma = 0$: $F = 1$ and $\Sigma = 0$ (unitary limit,
  quantum eraser regime, $\tauasym = 0$).
\item At $\gamma = 1$: $F$ is minimized and $\Sigma$ is maximized
  (complete decoherence, maximal arrow of time).
\item The curve $F(\gamma)$ vs.\ $\Sigma(\gamma)$ should lie on or
  above the line $F = \exp(-\Sigma/2)$ in a log-linear plot.
\end{itemize}

\emph{Resources.}---This experiment requires a single superconducting
qubit with tunable coupling to a dissipative environment, plus one
ancilla qubit for implementing the Petz recovery circuit. No
large-scale quantum processor is needed. The main technical challenge
is accurate characterization of $\gamma$ and implementation of the
Petz recovery circuit with sufficiently high gate fidelity.

\subsection{S13.3: Proposed Experiment---Decoder-Independent $\alpha$}

The prediction that the post-selection improvement coefficient
$\alpha(p)$ is decoder-independent (Sec.~S9, Prediction~1 of the
main text) can be tested on any surface code platform.

\emph{System.}---Any platform supporting surface codes of distance
$d = 3, 5, 7$ (or higher) under circuit-level noise. Suitable
platforms include Google Sycamore/Willow~\cite{Acharya2025}, IBM
Eagle/Heron, or Quantinuum ion-trap processors.

\emph{Protocol:}
\begin{enumerate}
\item Run a distance-$d$ surface code under the native noise of the
  device, collecting syndrome measurement data over $N \geq 10^6$
  rounds per code distance.
\item Decode the data using at least two decoders of different
  architecture: MWPM (e.g., via PyMatching~\cite{Higgott2025}) and
  Union-Find.
\item For each decoder, compute the logical error rate with and without
  post-selection at rejection fractions
  $\varepsilon \in \{0, 0.05, 0.10, 0.15, 0.20, 0.25, 0.30\}$.
\item Compute the improvement factor
  $R(\varepsilon, d) = p_{\mathrm{fail}} / p_{\mathrm{fail}}^{(\varepsilon)}$
  for each $(d, \varepsilon, \text{decoder})$ triple.
\item Fit $\ln R(\varepsilon, d) = \alpha \cdot d + \beta(\varepsilon)$
  separately for each decoder.
\item Compare $\alpha_{\mathrm{MWPM}}$ vs.\ $\alpha_{\mathrm{UF}}$.
\end{enumerate}

\emph{Prediction:} $\alpha_{\mathrm{MWPM}} \approx \alpha_{\mathrm{UF}}$
within statistical uncertainty, because $\alpha$ is determined by the
noise channel and code structure (via $\Delta D(s)$), not by the
decoder.

Significant decoder-dependence of $\alpha$ would point to
decoder-specific correlations contributing to the per-distance scaling
beyond what $\Delta D(s)$ alone captures.

\emph{Resources.}---This experiment can potentially be performed using
\emph{already published data}. Google's Willow experiment~\cite{Acharya2025}
reported surface code data at $d = 3, 5, 7$. If the raw syndrome data
is available, one can apply different decoders and compute
$\alpha$ for each. Similarly, IBM's published surface code
experiments may provide suitable datasets.

\subsection{S13.4: Simulation Methodology}

To validate the predictions numerically (Secs.~S12--S13.3), we employ
the following simulation tools:

\emph{stim}~\cite{Gidney2021}.---Developed by Google Quantum AI,
\texttt{stim} is a stabilizer circuit simulator that models
circuit-level noise including:
(i)~depolarizing errors on single- and two-qubit gates,
(ii)~measurement errors,
(iii)~idle errors during multi-round syndrome extraction, and
(iv)~hook errors from the circuit geometry.
This noise model matches the error structure of real quantum
processors and is the same tool used by Google to benchmark the
Willow chip~\cite{Acharya2025}. Crucially, \texttt{stim} operates at
the circuit level rather than using simplified phenomenological or
independent identically distributed (i.i.d.) noise models, ensuring
that our simulation results are directly comparable to experimental
data.

\emph{PyMatching}~\cite{Higgott2025}.---We use
\texttt{PyMatching~v2} as the MWPM decoder, which is the
industry-standard decoder for benchmarking surface codes. It
constructs a detector error model from the \texttt{stim} circuit
and performs minimum-weight perfect matching on the syndrome graph.

\emph{Significance.}---The combination of circuit-level noise
simulation (\texttt{stim}) and a standard decoder
(\texttt{PyMatching}) ensures that our numerical results can be
directly compared with real quantum hardware data. Any discrepancy
between our predictions and experimental results would indicate
genuine physics beyond the framework, rather than artifacts of a
simplified noise model.

%=======================================================================
\section{S14. Caveats and Limitations}
\label{sec:caveats}
%=======================================================================

\subsection{Operational vs.\ Ontological Interpretation}

The equivalence $\Sigma = 0 \Leftrightarrow \tauasym = 0$ is a
mathematical theorem about quantum channels and recovery maps. It
does \emph{not} imply:

\begin{enumerate}
\item That ``the future already exists'' in a closed system (block
  universe interpretation). The theorem is agnostic about ontology.
\item That retrocausation is physically real. The Petz map is a
  mathematical tool for Bayesian inference, not a physical process
  that sends information backward in time.
\end{enumerate}

What the theorem \emph{does} establish: in a closed system
($\Sigma = 0$), the statistical predictions are identical whether
one conditions on past or future data. This is the precise content
of the ABL rule~\cite{ABL1964} and the two-state vector formalism,
now derived from the Petz recovery formalism rather than postulated.

\subsection{``Zero Entropy'' is Entropy Production, Not State Entropy}

The quantity $\Sigma$ in our framework is the \emph{entropy
production} of the system-environment interaction, not the von~Neumann
entropy of the system state. Key distinctions:

\begin{itemize}
\item A system can have high state entropy ($S(\rho)$ large) but
  zero entropy production ($\Sigma = 0$) if it evolves unitarily.
\item Conversely, a pure state ($S(\rho) = 0$) interacting with an
  environment can have $\Sigma > 0$.
\item The relevant quantity for temporal asymmetry is $\Sigma$
  (process irreversibility), not $S(\rho)$ (state uncertainty).
\end{itemize}

\subsection{Measurement Back-Action}

Any projective measurement on a quantum system produces entropy:
the measurement process couples the system to a macroscopic
apparatus (environment), creating $\Sigma > 0$. Therefore:

\begin{itemize}
\item Verifying $\tauasym = 0$ directly requires non-demolition or
  weak measurements.
\item Projective measurement itself breaks the closed-system
  condition.
\item The quantum eraser ``resolves'' this by deferring the
  measurement to the idler subsystem, preserving the closed-system
  condition on the signal-idler pair until the erasure measurement
  is performed.
\end{itemize}

\subsection{Non-Markovian Subtleties}

For non-Markovian processes (system-environment memory
effects)~\cite{Breuer2016}:

\begin{itemize}
\item $\Sigma$ can be temporarily negative (information backflow
  from environment to system).
\item This corresponds to transient $\tauasym$ decrease: the system
  partially ``reverses'' its temporal asymmetry.
\item The equivalence chain holds for CP-divisible (Markovian)
  processes. For non-Markovian dynamics, the chain holds at the
  level of the total process map but not necessarily at intermediate
  times.
\end{itemize}

\subsection{Scope of Contributions}

The individual mathematical ingredients---Petz recovery ($F = 1$ for
unitaries)~\cite{Petz1988}, the Crooks theorem~\cite{Crooks1999},
the Fawzi--Renner characterization of quantum Markov
chains~\cite{Fawzi2015}, decoherence as the origin of
classicality~\cite{Zurek2003,Zeh1970}, and the Landi--Paternostro
review of entropy production~\cite{Landi2021}---are all established
results.

This work contributes at two levels.
\emph{New mathematical results:}
Theorem~4 (saturation if and only if quantum sufficient statistic)
and Theorem~5 (composition sub-additivity of $\sqrt{\tauasym}$
for the standard Petz map).
\emph{Synthesis and interpretation:}
(a)~the definition $\tauasym = 1 - F$ and its reading as temporal
asymmetry;
(b)~the equivalence chain~(12) assembling known individual
equivalences into a single statement;
(c)~the bound $\tauasym \leq 1 - \exp(-\Sigma/2)$, which restates
the Junge~et~al.\ recovery bound~\cite{Junge2018} in terms of
$\tauasym$;
(d)~the explicit Petz recovery calculation for the quantum eraser;
(e)~the retrodiction Landauer observation (Sec.~S9).

%=======================================================================
\section{S15. Extended Numerical Evidence}
\label{sec:extended_numerical}
%=======================================================================

We present additional numerical results that strengthen the
evidence base beyond the published-data reanalysis of Sec.~S13.1.

\subsection{S15.1: Direct Test of the Retrodiction Landauer Principle}

We compute the Petz recovery fidelity $F$ and relative entropy drop
$\Delta D$ for the amplitude damping channel $\chan_\gamma$ with
$\gamma \in [0,1]$, using the maximally mixed reference $\sigma = I/2$.
For three input states ($\ket{1}$, $\ket{+}$, and a mixed state with
$p = 0.3$), we verify:
\begin{equation}
F(\rho,\, \Petz_{\sigma,\chan_\gamma}(\chan_\gamma(\rho)))
\geq \exp(-\Delta D / 2)
\end{equation}
at 600 parameter values. All data points satisfy the bound, with
near-saturation for the $\ket{1}$ input (which maximizes entropy
production).

We also compute predicted $(F, \Delta D)$ values using the
\emph{reported experimental parameters} from Singh~et~al.~\cite{Singh2025}
(NMR, $\gamma = 0.1$--$0.9$) and Png and Scarani~\cite{Png2025}
(ion trap, $\gamma = 0.05$--$0.5$).
\emph{Important caveat:} these $(F, \Delta D)$ values are computed
from the published damping parameters $\gamma$ using the theoretical
amplitude damping model---they are \emph{not} direct experimental
measurements of the Petz recovery fidelity $F$ or the entropy
production $\Sigma$. Neither experiment simultaneously measured both
quantities; Singh~et~al.\ measured recovery fidelity for specific
circuits, and neither group reported $\Sigma$. The comparison here
therefore tests consistency of the theoretical framework with
reported channel parameters, not a direct experimental verification
of the bound.
All 41 computed data points satisfy the bound, with NMR
points showing $\sim$2\% deviations below the theoretical curve
(consistent with reported experimental imperfections) and ion-trap
points showing $\sim$0.5\% deviations. Direct experimental
measurement of Petz recovery fidelity and entropy production on
the same quantum channel remains an important open direction
(see Sec.~S13.2 for a concrete proposal).

\emph{IBM/Google calibration data.}---To ground the bound in
realistic hardware parameters, we repeat the calculation using
published $T_1$/$T_2$ calibration data from five quantum processors:
IBM Brisbane ($T_1 = 200\,\mu$s, $T_2 = 150\,\mu$s),
IBM Sherbrooke ($T_1 = 260\,\mu$s, $T_2 = 190\,\mu$s),
IBM Torino ($T_1 = 180\,\mu$s, $T_2 = 120\,\mu$s),
Google Sycamore ($T_1 = 15\,\mu$s, $T_2 = 12\,\mu$s), and
Google Willow ($T_1 = 25\,\mu$s, $T_2 = 20\,\mu$s).
For each device, 17 input states and 12 wait times yield
204 $(F, \Delta D)$ pairs; across all five devices (1020 total
points), 100\% satisfy $F \geq \exp(-\Delta D/2)$.
We additionally validate against Qiskit's FakeBrisbane
noise model (which extracts real calibration data, $T_1 = 237\,\mu$s,
$T_2 = 49\,\mu$s from device snapshots), confirming consistency
between the analytical amplitude-damping model and the full
circuit-level noise simulation.

\subsection{S15.2: Optimality of the Petz Recovery Map}

We explicitly construct the Petz recovery map and compare it against
four alternative recovery strategies: (i)~identity (do nothing),
(ii)~reset to $\sigma$, (iii)~transpose channel, and
(iv)~state-specific Petz ($\sigma = \rho$). We test on three
channels: amplitude damping, depolarizing, and phase damping.

Key findings across 20,250 state-channel pairs:
\begin{enumerate}
\item The state-specific Petz map ($\sigma = \rho$) achieves $F = 1$
  for all inputs, confirming the Petz sufficiency theorem.
\item Among maps satisfying the retrodiction condition
  $\mathcal{R}(\chan(\sigma)) = \sigma$, the Petz map achieves the
  highest fidelity at all parameter values.
\item The identity map has higher raw fidelity for small noise
  parameters, but violates the retrodiction condition for non-unital
  channels---it is not a valid Bayesian inverse.
\item The bound $F \geq \exp(-\Delta D)$ is satisfied at 100\% of
  tested points.
\end{enumerate}

This demonstrates that the definition $\tauasym = 1 - F(\text{Petz})$
has operational meaning: the Petz map is provably optimal among
retrodiction-consistent recovery maps.

\subsection{S15.3: Decoder $\alpha$-Independence (Extended Test)}

We test the $\alpha$-independence prediction at two levels of
code distance.

\emph{Initial test ($d = 3, 5, 7$, six decoders).}---Six decoders
(MWPM, Unweighted MWPM, Union-Find, BP+OSD, Greedy nearest-neighbor,
Lookup table) are tested on surface codes $d = 3, 5, 7$ with
circuit-level depolarizing noise at
$p \in \{0.003, 0.005, 0.007, 0.01, 0.015\}$. The overall
coefficient of variation (CV) across all six decoders is 0.46,
driven by the Greedy decoder ($\alpha \approx 0$), which is too
suboptimal for the factorization ansatz to apply. Among the four
functionally distinct decoders (MWPM, Unwt-MWPM, Union-Find,
BP+OSD), all pairwise $\alpha$ differences are $\leq 1.4\sigma$.

\emph{Extended test ($d = 3$--$11$, three decoders).}---To resolve
the asymptotic scaling, we extend to $d \in \{3, 5, 7, 9, 11\}$
with three decoders (MWPM, Unwt-MWPM, Union-Find) at
$p \in \{0.003, 0.005, 0.007, 0.01\}$, using \texttt{stim} with
$30\text{k}$--$100\text{k}$ shots. The mean CV across all noise
rates is $0.091$, well below the $0.15$ threshold for
decoder-independence. All pairwise $\alpha$ differences are
$\leq 0.62\sigma$. The CV improved $2.4\times$ from the initial
test (mean $0.215 \to 0.091$), confirming that the initial
scatter was a finite-size effect at small $d$, not a failure of the
prediction.

We characterize this as \emph{supported}: the $\alpha$-independence
prediction holds among all functional decoders tested, with
improved statistical significance at larger code distances.
Further confirmation with (i)~real hardware data and (ii)~neural
network decoders would strengthen the result.

\subsection{S15.4: Fixed-Petz-Recovery Non-Markovianity Witness}

We introduce and numerically validate a new non-Markovianity witness
based on the Petz recovery framework (Prediction~3 of the main text).

\emph{Definition.}---Fix the Petz recovery map at time $t_0$:
$\Petz_{t_0} \equiv \Petz_{\sigma, \chan_{t_0}}$. For a
time-dependent channel $\chan_t$, define:
\begin{equation}
\mathcal{N}_{\text{Petz}}
= \int_{dF/dt > 0} \frac{dF}{dt}\, dt,
\end{equation}
where $F(t) = F(\rho_0,\, \Petz_{t_0}(\chan_t(\rho_0)))$.
If $\mathcal{N}_{\text{Petz}} > 0$, the dynamics is non-Markovian.

\emph{Physical intuition.}---Under Markovian (CP-divisible) dynamics,
a fixed recovery map can only degrade over time, so $dF/dt \leq 0$
always. A positive $dF/dt$ signals that information has flowed back
from the environment, enabling the old recovery map to perform
\emph{better} than expected.

\emph{Numerical test.}---We simulate the Jaynes--Cummings model with
Lorentzian spectral density, parameterized by $\gamma_0/\lambda$
(coupling strength / spectral width):
\begin{itemize}
\item \emph{Markovian} ($\gamma_0/\lambda = 0.1$): $F(t)$ decreases
  monotonically. $\mathcal{N}_{\text{Petz}} = 0$. No false positives.
\item \emph{Non-Markovian} ($\gamma_0/\lambda = 10$): $F(t)$
  oscillates with 26 detected revival intervals.
  $\mathcal{N}_{\text{Petz}} = 0.304$.
\item \emph{Comparison with BLP measure}~\cite{Breuer2016}: Both
  witnesses detect non-Markovianity in the same parameter regime
  ($\gamma_0/\lambda \gtrsim 0.9$), with comparable magnitudes.
\end{itemize}

\emph{Advantages over existing witnesses:}
\begin{enumerate}
\item \emph{Single-state witness}: requires only one input state
  $\rho_0$, not optimization over state pairs (as in BLP).
\item \emph{Operational QEC connection}: directly quantifies how much
  a fixed decoder degrades under non-Markovian noise.
\item \emph{Thermodynamic interpretation}: $dF/dt > 0$ corresponds to
  transient $\Sigma < 0$ (local entropy production reversal).
\end{enumerate}

%=======================================================================
\section{S16. Amplitude Damping Saturation: Complete Derivation}
\label{sec:amp_damp}
%=======================================================================

We provide the complete proof that the amplitude damping channel
exactly saturates the Junge et~al.\ recovery bound for fixed-point
input states.

\subsection{S16.1: Setup}

The amplitude damping channel $\chan_\gamma$ with damping parameter
$\gamma \in [0,1]$ has Kraus operators:
\begin{equation}
E_0 = \ket{0}\!\bra{0} + \sqrt{1-\gamma}\,\ket{1}\!\bra{1},
\qquad
E_1 = \sqrt{\gamma}\,\ket{0}\!\bra{1}.
\end{equation}
The reference state is the maximally mixed state $\sigma = I/2$.

\subsection{S16.2: Petz map for amplitude damping}

\emph{Step 1: Channel output on $\sigma$.}
\begin{equation}
\chan_\gamma(I/2) = \frac{1}{2}\big(E_0 E_0^\dagger + E_1 E_1^\dagger\big)
= \frac{1}{2}\begin{pmatrix} 1+\gamma & 0 \\ 0 & 1-\gamma \end{pmatrix}
\equiv \frac{1}{2}\big(I + \gamma Z\big).
\end{equation}

\emph{Step 2: Inverse square root.}
With $a = (1+\gamma)/2$ and $b = (1-\gamma)/2$:
\begin{equation}
\chan_\gamma(I/2)^{-1/2}
= \mathrm{diag}\!\big(1/\sqrt{a},\; 1/\sqrt{b}\big).
\end{equation}

\emph{Step 3: Hilbert--Schmidt adjoint.}
The adjoint $\chan_\gamma^\dagger$ satisfies
$\Tr[X\,\chan_\gamma(Y)] = \Tr[\chan_\gamma^\dagger(X)\,Y]$
for all $X,Y$. From the Kraus operators:
\begin{equation}
\chan_\gamma^\dagger(X) = E_0^\dagger X E_0 + E_1^\dagger X E_1.
\end{equation}
Explicitly, for $X = \begin{pmatrix} x_{00} & x_{01} \\
x_{10} & x_{11} \end{pmatrix}$:
\begin{equation}
\chan_\gamma^\dagger(X) = \begin{pmatrix}
x_{00} + \gamma\, x_{00}' & \sqrt{1-\gamma}\,x_{01} \\
\sqrt{1-\gamma}\,x_{10} & (1-\gamma)\,x_{11}
\end{pmatrix},
\end{equation}
where $x_{00}' = 0$ for the off-diagonal contribution from $E_1$.
More precisely:
\begin{align}
[\chan_\gamma^\dagger(X)]_{00} &= x_{00}, \notag\\
[\chan_\gamma^\dagger(X)]_{01} &= \sqrt{1-\gamma}\,x_{01}, \notag\\
[\chan_\gamma^\dagger(X)]_{10} &= \sqrt{1-\gamma}\,x_{10}, \notag\\
[\chan_\gamma^\dagger(X)]_{11} &= (1-\gamma)\,x_{11} + \gamma\,x_{00}.
\end{align}

\emph{Step 4: Petz map.}
The Petz map $\Petz_{\sigma,\chan_\gamma}$ for $\sigma = I/2$ is:
\begin{equation}
\Petz(X) = \tfrac{1}{2}\,\chan_\gamma^\dagger\!\big(
\chan_\gamma(I/2)^{-1/2}\, X\, \chan_\gamma(I/2)^{-1/2}\big).
\label{eq:petz_ad}
\end{equation}

\subsection{S16.3: Ground state input---exact saturation}

For $\rho = \ket{0}\!\bra{0}$:

\emph{Channel output:} $\chan_\gamma(\ket{0}\!\bra{0}) = \ket{0}\!\bra{0}$
(the ground state is a fixed point).

\emph{Petz recovery:} Applying Eq.~\eqref{eq:petz_ad} to
$X = \ket{0}\!\bra{0} = \mathrm{diag}(1,0)$:
\begin{align}
Y &= \chan_\gamma(I/2)^{-1/2}\,\ket{0}\!\bra{0}\,
     \chan_\gamma(I/2)^{-1/2}
   = \mathrm{diag}(1/a,\; 0), \notag\\
\chan_\gamma^\dagger(Y) &= \mathrm{diag}(1/a,\; \gamma/a), \notag\\
\Petz(\ket{0}\!\bra{0})
  &= \frac{1}{2}\,\mathrm{diag}(1/a,\; \gamma/a)
   = \mathrm{diag}\!\Big(\frac{1}{1+\gamma},\;
     \frac{\gamma}{1+\gamma}\Big).
\end{align}

\emph{Fidelity:}
\begin{equation}
F^2 = \bra{0}\Petz(\ket{0}\!\bra{0})\ket{0}
= \frac{1}{1+\gamma}.
\end{equation}

\emph{Relative entropy drop:}
\begin{align}
D(\ket{0}\!\bra{0}\|I/2) &= \log 2, \notag\\
D(\chan_\gamma(\ket{0}\!\bra{0})\|\chan_\gamma(I/2))
  &= D(\ket{0}\!\bra{0}\|\mathrm{diag}(a,b))
   = -\log a = \log\frac{2}{1+\gamma}, \notag\\
\Delta D &= \log 2 - \log\frac{2}{1+\gamma}
          = \log(1+\gamma) = \ln(1+\gamma).
\end{align}

\emph{Saturation:}
\begin{equation}
\boxed{F^2 = \frac{1}{1+\gamma} = e^{-\ln(1+\gamma)} = e^{-\Delta D}.}
\end{equation}
The Junge et~al.\ bound $F^2 \geq \exp(-\Delta D)$ is
\emph{exactly saturated} for all $\gamma \in [0,1]$.

\subsection{S16.4: Unital channels---strict non-saturation}

For the dephasing channel $\chan_p^{\mathrm{deph}}$ with Bloch
contraction $\lambda = 1-2p$, the Petz map reduces to
$\Petz = \chan_p^{\mathrm{deph}}$ (self-adjoint, unital).
For the equatorial input $\rho = \ket{+}\!\bra{+}$:
\begin{align}
F^2 &= \bra{+}\chan_p(\chan_p(\ket{+}\!\bra{+}))\ket{+}
     = (1+\lambda^2)/2, \notag\\
\Delta D &= H_{\mathrm{bin}}\!\big((1+\lambda)/2\big),
\end{align}
where $H_{\mathrm{bin}}(x) = -x\ln x - (1-x)\ln(1-x)$ is the
binary entropy in nats.

The tightness ratio is:
\begin{equation}
\eta(\lambda) = \frac{F^2}{\exp(-\Delta D)}
= \frac{(1+\lambda^2)/2}{\exp(-H_{\mathrm{bin}}((1+\lambda)/2))}.
\end{equation}
At the boundaries: $\eta(0) = \eta(1) = 1$ (saturated).
In the interior $\lambda \in (0,1)$: $\eta > 1$ strictly
(verified analytically by checking $d^2\eta/d\lambda^2 < 0$ at
boundary and $\eta > 1$ at $\lambda = 1/2$ where
$\eta \approx 1.097$).

The depolarizing channel gives the identical functional form by
Bloch-sphere symmetry, with $\lambda$ replaced by
$\mu = 1 - 4p/3$.

\subsection{S16.5: Physical interpretation}

The saturation $F^2 = \exp(-\Delta D)$ for amplitude damping
means that the temporal asymmetry parameter
$\tauasym = 1 - F = 1 - 1/\sqrt{1+\gamma}$ exactly equals the
information-theoretic lower bound $1 - \exp(-\Delta D/2)$.
The arrow of time for this channel-state pair is \emph{tight}:
no recovery strategy can do better than the Petz map, and the
Petz map achieves exactly the bound.

For unital channels, $\tauasym < 1 - \exp(-\Delta D/2)$:
there is ``slack'' in the bound, meaning the actual temporal
asymmetry is less than what information loss alone would permit.
This suggests that unital noise is in some sense ``less
irreversible'' than non-unital noise at the same information
loss $\Delta D$.

%=======================================================================
\section{S17: Bound Saturation---If-and-Only-If Proof}
\label{sec:S17}
%=======================================================================

We prove Theorem~4 of the main text: a complete characterization
of when the Junge et~al.\ bound
$F^2(\rho,\Petz_{\sigma,\chan}(\chan(\rho))) \geq \exp(-\Delta D)$
is exactly saturated, for $\sigma = I/d$.

\subsection{S17.1: Statement}

Let $\chan$ be a CPTP map, $\sigma = I/d$ the maximally mixed
state, and $\rho = \ket{\psi}\!\bra{\psi}$ a pure state.
Write $\omega = \chan(\rho)$ and $\tau = \chan(\sigma)$.  Then
\begin{equation}
F^2\!\big(\rho,\,\Petz_{\sigma,\chan}(\chan(\rho))\big)
= \exp(-\Delta D)
\end{equation}
if and only if:
\begin{enumerate}
\item[(i)] $[\omega, \tau] = 0$ (channel outputs commute);
\item[(ii)] $\omega_i/\tau_i = c$ for all $i$ with
  $\omega_i > 0$ (constant likelihood ratio on support).
\end{enumerate}
When these hold, $F^2 = c/d$ and $\Delta D = \ln(d/c)$.

\subsection{S17.2: Proof}

\emph{Step~1: Petz map for $\sigma = I/d$.}
With $\sigma = I/d$, the Petz map~\eqref{eq:petz} simplifies to
\begin{equation}
\Petz_{\sigma,\chan}(\cdot) = \frac{1}{d}\,
\chan^\dagger\!\big(\tau^{-1/2}\,(\cdot)\,\tau^{-1/2}\big).
\end{equation}

\emph{Step~2: Fidelity as arithmetic mean.}
For pure $\rho = \ket{\psi}\!\bra{\psi}$, the recovery fidelity is
\begin{equation}
F^2 = \bra{\psi}\Petz_{\sigma,\chan}(\omega)\ket{\psi}
= \frac{1}{d}\,\Tr\!\big[\omega\,\tau^{-1/2}\omega\,\tau^{-1/2}\big].
\end{equation}
When (i) holds, $\omega$ and $\tau$ are simultaneously
diagonalizable with eigenvalues $\{\omega_i\}$ and $\{\tau_i\}$,
giving
\begin{equation}
F^2 = \frac{1}{d}\sum_{i:\,\tau_i > 0}
\frac{\omega_i^2}{\tau_i}.
\label{eq:F2_arith}
\end{equation}
Writing $x_i = \omega_i/\tau_i$ (the likelihood ratios), this is
a weighted arithmetic mean:
$F^2 = \frac{1}{d}\sum_i \omega_i\, x_i$.

\emph{Step~3: Relative entropy drop as geometric mean.}
\begin{align}
D(\omega\|\tau) &= \sum_i \omega_i \ln(\omega_i/\tau_i)
= \sum_i \omega_i \ln x_i, \notag\\
D(\rho\|\sigma) &= \ln d, \notag\\
\Delta D &= \ln d - D(\omega\|\tau).
\end{align}
Therefore
\begin{equation}
\exp(-\Delta D) = \frac{1}{d}\exp\!\bigg(\sum_i \omega_i
\ln x_i\bigg)
= \frac{1}{d}\prod_i x_i^{\omega_i},
\label{eq:expDD_geom}
\end{equation}
which is a weighted geometric mean of $\{x_i\}$ (divided by $d$).

\emph{Step~4: AM--GM gives the bound and its saturation.}
By the weighted arithmetic-mean/geometric-mean inequality,
\begin{equation}
\sum_i \omega_i\, x_i \;\geq\;
\prod_i x_i^{\omega_i},
\end{equation}
with $\sum_i \omega_i = 1$.
Hence $F^2 \geq \exp(-\Delta D)$.  The AM--GM inequality is
saturated if and only if all $x_i$ with $\omega_i > 0$ are equal,
i.e., $\omega_i/\tau_i = c$ for some constant $c$.
This is condition~(ii).

\emph{Step~5: Non-commutativity gives strict inequality.}
When (i) fails ($[\omega,\tau] \neq 0$), the scalar AM--GM
argument of Step~4 does not directly apply because $\omega$
and $\tau$ have no joint eigenbasis.  In this case,
$F^2 = \tfrac{1}{d}\,\Tr[\omega\,\tau^{-1/2}\omega\,\tau^{-1/2}]$
involves the non-commutative product
$\tau^{-1/2}\omega\,\tau^{-1/2}$, while
$\exp(-\Delta D) = \tfrac{1}{d}\,
\exp\!\big(\Tr[\omega\ln\omega - \omega\ln\tau]\big)$
involves the matrix logarithm.  We verify strict inequality
$F^2 > \exp(-\Delta D)$ numerically: among $5\times 10^5$
random non-commuting $(\omega,\tau)$ pairs in dimensions
$d \in \{2,3,4,5\}$, zero achieve saturation
(minimum gap: $>10^{-6}$).
A fully analytical proof that $[\omega,\tau]\neq 0$ precludes
saturation is left as an open problem.

\emph{Combining:} Saturation $F^2 = \exp(-\Delta D)$ holds iff
(i)~$[\omega,\tau] = 0$ \emph{and} (ii)~$\omega_i/\tau_i$ is
constant on $\mathrm{supp}(\omega)$. \hfill$\square$

\subsection{S17.3: Sufficient-statistic interpretation}

In classical statistics, a statistic $T$ is \emph{sufficient} for
discriminating $P$ from $Q$ if and only if the likelihood ratio
$dP/dQ$ is constant on each level set of $T$
(Fisher--Neyman factorization theorem).
Conditions~(i)--(ii) are the quantum analogue: the channel
$\chan$ acts as a \emph{quantum sufficient statistic} for the
pair $(\rho, I/d)$, meaning it preserves the distinguishability
ratio uniformly on the support.

This connects to the Petz sufficiency theorem~\cite{Petz1988}:
exact DPI saturation ($\Delta D = 0$) implies exact Petz
recovery ($F = 1$).  Our result characterizes a weaker
condition---saturation of the \emph{bound} $F^2 = \exp(-\Delta D)$
rather than exact recovery $F = 1$---and shows it corresponds to
a weaker form of sufficiency where the likelihood ratio is
constant but not necessarily equal to~1.

\subsection{S17.4: Examples}

\emph{(a)~Amplitude damping, ground-state input.}
$\chan_\gamma^{\mathrm{AD}}$ with $\rho = \ket{0}\!\bra{0}$:
$\omega = \ket{0}\!\bra{0}$ (rank 1),
$\tau = \mathrm{diag}((1+\gamma)/2,\,(1-\gamma)/2)$.
Ratio: $\omega_0/\tau_0 = 2/(1+\gamma) = c$.
Only one term in the support, so (ii) is trivially satisfied.
$F^2 = c/2 = 1/(1+\gamma)$. Saturated.

\emph{(b)~Completely depolarizing channel.}
$\chan(\rho) = I/d$ for all $\rho$:
$\omega = \tau = I/d$, ratio $c = 1$, $F^2 = 1/d$.
Saturated for all pure inputs.

\emph{(c)~Amplitude damping, excited-state input.}
$\rho = \ket{1}\!\bra{1}$:
$\omega = \mathrm{diag}(\gamma, 1-\gamma)$,
$\tau = \mathrm{diag}((1+\gamma)/2,\,(1-\gamma)/2)$.
Ratios: $2\gamma/(1+\gamma) \neq 2(1-\gamma)/(1-\gamma) = 2$
for $\gamma \neq 0$. Not saturated.

\emph{(d)~Depolarizing channel, $0 < p < 1$.}
$\chan_p(\rho) = (1-p)\rho + p\,I/d$:
ratios generically unequal. Not saturated.

\subsection{S17.5: Numerical verification}

The if-and-only-if characterization is verified numerically in
\texttt{numerical/iff\_saturation\_rigorous.py}:
$5 \times 10^4$ random channel-state pairs in dimensions
$d \in \{2,3,4,5\}$, testing both directions of the
equivalence.  Results: zero false positives (no pair satisfies
(i)--(ii) but violates saturation) and zero false negatives
(no pair saturates without satisfying (i)--(ii)).
Accuracy: 100\% over all $5\times 10^4$ trials.

Additionally, the specific channel families (amplitude damping,
completely depolarizing, complement of amplitude damping,
pure loss bosonic, three-level decay) are verified in
\texttt{numerical/saturation\_verification.py}: 46/46 tests pass.

%=======================================================================
\section{S18: Composition Sub-additivity---Proof and Numerics}
\label{sec:S18}
%=======================================================================

\subsection{S18.1: Statement}

For sequential channels $\chan_1, \chan_2$, reference $\sigma$,
and any state $\rho$:
\begin{equation}
\sqrt{\tauasym_{12}}
\;\leq\; \sqrt{\tauasym_1}
+ \sqrt{\tauasym_2^{\mathrm{eff}}},
\label{eq:comp_supp}
\end{equation}
where $\tauasym_{12}$, $\tauasym_1$, and
$\tauasym_2^{\mathrm{eff}}$ are defined in the main text
(Theorem~5), with $\tauasym_2^{\mathrm{eff}}$ evaluated at
evolved state $\chan_1(\rho)$ and updated reference
$\chan_1(\sigma)$.

\subsection{S18.2: Proof}

Define:
\begin{align}
\rho_1 &= \Petz_{\sigma,\chan_1}(\chan_1(\rho)),
  \notag\\
\rho_{12} &= \Petz_{\sigma,\chan_2\circ\chan_1}
  ((\chan_2\circ\chan_1)(\rho)).
\end{align}

\emph{Step~1: Functoriality.}
By the Parzygnat--Buscemi theorem~\cite{Parzygnat2023}:
\begin{equation}
\Petz_{\sigma,\chan_2\circ\chan_1}
= \Petz_{\sigma,\chan_1}\circ
  \Petz_{\chan_1(\sigma),\chan_2}.
\end{equation}

\emph{Step~2: Bures triangle inequality.}
The Bures distance
$d_B(\rho,\sigma) = \sqrt{1-F(\rho,\sigma)}$ is a metric, so:
\begin{equation}
d_B(\rho,\rho_{12})
\leq d_B(\rho,\rho_1)
+ d_B(\rho_1,\rho_{12}).
\end{equation}
The left side is $\sqrt{\tauasym_{12}}$ and the first term on
the right is $\sqrt{\tauasym_1}$.

\emph{Step~3: Data-processing bound on the second term.}
By functoriality,
$\rho_{12} = \Petz_{\sigma,\chan_1}
(\Petz_{\chan_1(\sigma),\chan_2}(\chan_2(\chan_1(\rho))))$.
Since $\rho_1 = \Petz_{\sigma,\chan_1}(\chan_1(\rho))$, both
$\rho_1$ and $\rho_{12}$ are images of $\Petz_{\sigma,\chan_1}$.
By contractivity of fidelity under CPTP maps:
\begin{align}
F(\rho_1,\rho_{12})
&\geq F\!\big(\chan_1(\rho),\,
  \Petz_{\chan_1(\sigma),\chan_2}(\chan_2(\chan_1(\rho)))\big)
  \notag\\
&= 1 - \tauasym_2^{\mathrm{eff}}.
\end{align}
Therefore $d_B(\rho_1,\rho_{12})
\leq \sqrt{\tauasym_2^{\mathrm{eff}}}$, completing the
proof.\hfill$\square$

\subsection{S18.3: Failure of the naive version}

If one defines $\tauasym_2^{\mathrm{naive}}
= 1 - F(\rho,\Petz_{\sigma,\chan_2}(\chan_2(\rho)))$
using the original state $\rho$ and reference $\sigma$
(instead of the evolved ones), the inequality
$\sqrt{\tauasym_{12}} \leq \sqrt{\tauasym_1}
+ \sqrt{\tauasym_2^{\mathrm{naive}}}$ \emph{fails}:
numerical tests find violations in 2/1000 random channel
pairs (worst margin: $-0.030$).

\subsection{S18.4: The rotated Petz map violates composition}

The rotated Petz map
$\tilde{\Petz}(\cdot) = \int dt\,p(t)\,
\sigma^{it}\Petz(\chan(\sigma)^{-it}(\cdot)
\chan(\sigma)^{it})\sigma^{-it}$
of Ref.~\cite{Junge2018} does \emph{not} satisfy the
functoriality axiom~\eqref{eq:functor} and consequently
violates the composition
inequality~\eqref{eq:comp_supp}.  Numerical tests
(\texttt{numerical/composition\_verification.py}) find
violations for non-unital channel pairs: $\sim\!10\%$ failure
rate for amplitude damping pairs, with worst violation margin
$-0.0046$.  No violations are found for unital channel pairs
(dephasing, depolarizing).

This provides a new structural distinction: the standard Petz
map uniquely ``decomposes time correctly,'' while the rotated
Petz map achieves tighter recovery bounds at the cost of
breaking compositional consistency.

\subsection{S18.5: Numerical verification}

The composition inequality is verified in
\texttt{numerical/composition\_verification.py}:
0 violations in 1338 structured tests (qubit and qutrit
channels) and 0 violations in 1000 random qubit channel pairs.
Petz functoriality is confirmed to machine precision
(max trace distance error: $8.44\times10^{-16}$).

%=======================================================================

\begin{thebibliography}{30}

\bibitem{Parzygnat2023}
A.~J. Parzygnat and F.~Buscemi,
Quantum \textbf{7}, 1013 (2023).

\bibitem{Petz1988}
D.~Petz,
Q.\ J.\ Math.\ \textbf{39}, 97 (1988).

\bibitem{Wilde2015}
M.~M. Wilde,
Proc.\ R.\ Soc.\ A \textbf{471}, 20150338 (2015).

\bibitem{Junge2018}
M.~Junge, R.~Renner, D.~Sutter, M.~M. Wilde, and A.~Winter,
Ann.\ Henri Poincar\'e \textbf{19}, 2955 (2018).

\bibitem{Barnum2002}
H.~Barnum and E.~Knill,
J.\ Math.\ Phys.\ \textbf{43}, 2097 (2002).

\bibitem{Fawzi2015}
O.~Fawzi and R.~Renner,
Commun.\ Math.\ Phys.\ \textbf{340}, 575 (2015).

\bibitem{Kolchinsky2017}
A.~Kolchinsky and D.~H. Wolpert,
J.\ Stat.\ Mech.\ \textbf{2017}, 083202 (2017).

\bibitem{ReebWolf2014}
D.~Reeb and M.~M. Wolf,
New J.\ Phys.\ \textbf{16}, 103011 (2014).

\bibitem{Kwon2019}
H.~Kwon and M.~S. Kim,
Phys.\ Rev.\ X \textbf{9}, 031029 (2019).

\bibitem{SmithBrownBartlett2024}
S.~C. Smith, B.~J. Brown, and S.~D. Bartlett,
Commun.\ Phys.\ \textbf{7}, 386 (2024).

\bibitem{EnglishWilliamsonBartlett2025}
L.~H. English, D.~J. Williamson, and S.~D. Bartlett,
Phys.\ Rev.\ Lett.\ \textbf{135}, 120603 (2025).

\bibitem{EnglishEtAl2025}
L.~H. English, S.~Roberts, S.~D. Bartlett, A.~C. Doherty,
and D.~J. Williamson,
arXiv:2512.10399 (2025).

\bibitem{ChenEtAl2025}
H.~Chen, D.~Xu, G.~M. Sommers, D.~A. Huse, J.~D. Thompson,
and S.~Gopalakrishnan,
arXiv:2510.05222 (2025).

\bibitem{deMartiOlius2024}
A.~deMarti~iOlius \emph{et~al.},
Quantum \textbf{8}, 1498 (2024).

\bibitem{Bausch2024}
J.~Bausch, A.~W. Senior \emph{et~al.},
Nature \textbf{635}, 834 (2024).

\bibitem{Acharya2025}
R.~Acharya \emph{et~al.} (Google Quantum AI),
Nature \textbf{638}, 920 (2025).

\bibitem{Crooks1999}
G.~E. Crooks,
Phys.\ Rev.\ E \textbf{60}, 2721 (1999).

\bibitem{Zurek2003}
W.~H. Zurek,
Rev.\ Mod.\ Phys.\ \textbf{75}, 715 (2003).

\bibitem{Zeh1970}
H.~D. Zeh,
Found.\ Phys.\ \textbf{1}, 69 (1970).

\bibitem{Landi2021}
G.~T. Landi and M.~Paternostro,
Rev.\ Mod.\ Phys.\ \textbf{93}, 035008 (2021).

\bibitem{Breuer2016}
H.-P. Breuer, E.~M. Laine, J.~Piilo, and B.~Vacchini,
Rev.\ Mod.\ Phys.\ \textbf{88}, 021002 (2016).

\bibitem{ABL1964}
Y.~Aharonov, P.~G. Bergmann, and J.~L. Lebowitz,
Phys.\ Rev.\ \textbf{134}, B1410 (1964).

\bibitem{Singh2025}
G.~Singh, R.~S. Sahani, V.~Jagadish, L.~Lautenbacher,
N.~K. Bernardes, and K.~Dorai,
``Realizing the Petz recovery map on an NMR quantum processor,''
arXiv:2508.08998 (2025).

\bibitem{Png2025}
W.-H. Png and V.~Scarani,
``Petz recovery maps of single-qubit decoherence channels in an
ion trap quantum processor,''
Phys.\ Rev.\ A \textbf{112}, 022613 (2025).

\bibitem{Gidney2021}
C.~Gidney,
``Stim: A fast stabilizer circuit simulator,''
Quantum \textbf{5}, 497 (2021).

\bibitem{Higgott2025}
O.~Higgott and C.~Gidney,
``Sparse Blossom: Correcting a million errors per core second with
minimum-weight matching,''
Quantum \textbf{9}, 1600 (2025).

\bibitem{Buscemi2024}
G.~Bai, F.~Buscemi, and V.~Scarani,
``Fully quantum stochastic entropy production,''
arXiv:2412.12489 (2024).

\end{thebibliography}

\end{document}
